\documentclass[oneside, 11pt]{amsbook}
\linespread{1.1}

\def\a{\alpha}
\def\Id{\rm Id}
%\UseRawInputEncoding
\usepackage[margin=1in]{geometry}
\usepackage{graphics}
\usepackage[english]{babel}
\usepackage{amssymb,amsmath,amsfonts}
\usepackage{datetime}
\usepackage{color}
\usepackage{ulem}
\usepackage{lipsum}
\usepackage{amsfonts}
\usepackage[utf8]{inputenc}
\usepackage{cite}
\synctex=1
\RequirePackage{amscd}
\usepackage{marginnote}
\newcommand{\margem}[1]{\marginpar{{\scriptsize {#1}}}}

\RequirePackage{epic}
\RequirePackage[all]{xy}
\RequirePackage{url}
\RequirePackage[shortlabels]{enumitem}
\usepackage{empheq}
\usepackage{epsfig}
\usepackage{hyperref}
\usepackage{cleveref}
\usepackage[english]{babel}

\usepackage{draftwatermark}
\SetWatermarkText{Draft: \today}
\SetWatermarkColor[gray]{0.5}
\SetWatermarkFontSize{1cm}
\SetWatermarkAngle{90}
\SetWatermarkHorCenter{20cm}
%\newcommand{\margem}[1]{\marginpar{{\scriptsize {#1}}}}
\usepackage[utf8]{inputenc}
\usepackage{cite}
%\usepackage{geometry}
\synctex=1
\RequirePackage{amscd}
%\RequirePackage[dvips]{graphicx}
\RequirePackage{epic}
\RequirePackage{eepic}
\RequirePackage[all]{xy}
\RequirePackage{url}
\RequirePackage[shortlabels]{enumitem}
\usepackage{empheq}
\usepackage{nomencl}
\usepackage{epsfig}
\usepackage{graphicx}
%%%%%%%%%%%%%%%%%%%%%%%%%%%%%%%%%% macros %%%%%%%%%%%%%%%%%%%%%%%%%%%%%%%%%%%%%%%%%

\newcommand{\comment}[1]{}
\renewcommand{\labelitemi}{$-$}   
    \newcommand{\set}[1]{{\left\{#1\right\}}}
\newcommand{\pa}[1]{{\left(#1\right)}}
\newcommand{\sq}[1]{{\left[#1\right]}}
\newcommand{\gen}[1]{{\left\langle #1\right\rangle}}
\newcommand{\abs}[1]{{\left|#1\right|}}
\newcommand{\norm}[1]{{\left\|#1\right\|}}
\newcommand{\res}[1]{\left.#1\right|} %restriction
\newcommand{\normal}{\mathrm}
\newcommand{\G}{\mathcal{G}}
\newcommand{\A}{\mathcal{A}}
 \newcommand{\U}{\mathcal{U}}
\newcommand{\V}{\mathcal{V}}
\newcommand{\qpsys}{(q,p)=(q_1,\ldots,q_n,p_1,\ldots,p_n)}
\newcommand{\T}{\mathbb{T}}
\newcommand{\Z}{\mathbb{Z}}
\newcommand{\R}{\mathbb{R}}
\newcommand{\C}{\mathbb{C}}
\newcommand{\teta}{\theta}
\newcommand{\eps}{\varepsilon}
\renewcommand{\Re}{\operatorname{Re}}
\renewcommand{\Im}{\operatorname{Im}}
\newcommand{\sgn}{\operatorname{sgn}}
\newcommand{\tr}{\operatorname{tr}}
\newcommand{\Mat}{\operatorname{Mat}}
\newcommand{\GL}{\operatorname{GL}}

\newcommand{\mat}[1]{{\begin{matrix}#1\end{matrix}}}
\newcommand{\pmat}[1]{{\begin{pmatrix}#1\end{pmatrix}}}
\newcommand{\amat}[2]{{\begin{array}{#1}#2\end{array}}}
\newcommand{\eqsys}[1]{{\left\{\begin{array}{ll}#1\end{array}\right.}}
\newcommand{\conj}{h\circ R_{2\pi\alpha}\circ h^{-1}}
\newcommand{\Ham}{\operatorname{Ham}}
\newcommand{\h}{\operatorname{H}}
\newcommand{\average}[1]{\int_{\T^n}#1}
\newcommand{\media}[1]{\int_\T #1}
\newcommand{\uh}{u^{\text{H}}} % Hamiltonian u
\newcommand{\vh}{v^{\text{H}}} % Hamiltonian v 
\newcommand{\Uh}{\mathcal{U}^{\operatorname{Ham}}}
\newcommand{\Vh}{\mathcal{V}^{\operatorname{Ham}}}
\newcommand{\Vhz}{\mathcal{V}^{\operatorname{Ham}}\oplus (-\eta r + \eta\Omega\,\base{r})}
\newcommand{\Vhd}{\mathcal{V}^{\operatorname{Ham}}\oplus (-\eta (r + \Omega)\base{r})}
\newcommand{\Uhd}{\mathcal{U}^{\operatorname{Ham}}\oplus (-\eta r\base{r})}
\newcommand{\tensor}{\otimes}% prodotto tensore
\newcommand{\inderiv}[1]{#1'^{-1}}    
\newcommand{\pull}[1]{#1^\ast} %pull back
\newcommand{\push}[1]{#1_\ast}%push forward
\newcommand{\base}[1]{{\frac{\partial}{\partial{#1}}}} %base campo vettoriale
\newcommand{\elh}[1]{{\frac{\partial #1}{\partial q}-\frac{d}{dt}\frac{\partial #1}{\partial\dot{q}}}}
\newcommand{\eli}[1]{{\frac{\partial #1}{\partial q_i}-\frac{d}{dt}\frac{\partial #1}{\partial\dot{q}_i}}}
\newcommand{\ph}[1]{\frac{\partial H}{\partial #1}}
\newcommand{\co}[1]{\textit{#1}}%corsivo
\newcommand{\gr}[1]{\textbf{#1}}%grassetto
\newcommand{\bla}[1]{\int_0^{{#1}}}
\newcommand{\tl}[1]{\tilde{#1}}
\newcommand{\blas}[1]{\int_{t_0}^{t_{#1}}}
\newcommand{\upp}[1]{\uppercase{#1}}
\newcommand{\id}{\operatorname{id}}
\newcommand{\Lip}{\operatorname{Lip}}
\newcommand{\Gr}{\operatorname{Gr}}
\newcommand{\heta}{{\eta}}
%%%%%%%%%%%%%%%%%%%%%% fine macros %%%%%%
% % % % % % % % % % % % % % % % % % % % %macromi % % % % % % % % % % % % % % % % % % % % % % % % % % % % %
\newcommand{\wtd}{{\widetilde{d}}}
\newcommand{\wtg}{{\widetilde{g}}}
\newcommand{\g}{\gamma}
\newcommand{\wth}{{\widetilde{h}}}
\newcommand{\wtp}{{\widetilde{p}}}
\newcommand{\wtq}{{\widetilde{q}}}
\newcommand{\wtr}{{\widetilde{r}}}
\newcommand{\wts}{{\widetilde{s}}}
\newcommand{\wtu}{{\widetilde{u}}}
\newcommand{\wtx}{{\widetilde{x}}}
\newcommand{\wtz}{{\widetilde{z}}}
\newcommand{\wty}{{\widetilde{y}}}
\newcommand{\wtw}{{\widetilde{w}}}
\newcommand{\new}{{\small{\mathtt{ New}}}}
\newcommand{\yy}{{\mathcal Y}}%{{ \mbox{\Fontlukas Y}}}
\newcommand{\s}{{\sigma}}
%\newcommand{\id}{{\rm Id}}
\newcommand{\jj}{{\vec \jmath}}
\newcommand{\ii}{{\vec \imath}}
\newcommand{\li}{{\mathfrak h}}
% widetilde uppercase letters
\newcommand{\wtH}{\widetilde{H}}
\newcommand{\wtK}{\widetilde{K}}
\newcommand{\wtL}{{\widetilde L}}
\newcommand{\wtM}{\widetilde{M}}
\newcommand{\wtN}{\widetilde{N}}
\newcommand{\wtP}{\widetilde{P}}
\newcommand{\wtQ}{\widetilde{Q}}
\newcommand{\wtR}{\widetilde{R}}
\newcommand{\wtS}{\widetilde{S}}
\newcommand{\wtT}{\widetilde{T}}
\newcommand{\wtU}{{\widetilde U}}
\newcommand{\wtV}{\widetilde{V}}
\newcommand{\wtX}{\widetilde{X}}
\newcommand{\wtZ}{\widetilde{Z}}
\def\wcK{\widetilde{\cK}}
\def\wwcK{\widetilde{\widetilde{\cK}}}
\def\norma#1{\left\|#1\right\|}
\def\sleq{\preceq}
%\newcommand{\ad}{{\rm ad}}
\newcommand{\ccC}{{\mathscr{C}}}
\newcommand{\ccL}{{\mathscr{L}}}
\newcommand{\om}{{\omega}}
% widetilde misc letters
\newcommand{\wt}[1]{{\widetilde{#1}}}
\newcommand{\wtomega}{\widetilde{\omega}}
\newcommand{\wtpi}{{\widetilde{\pi}}}
\newcommand{\wtnu}{{\widetilde{\nu}}}
\newcommand{\wtiota}{{\widetilde{\iota}}}
\newcommand{\wteta}{\ex{-s^j}\eta_{\ps}}
\newcommand{\wtchi}{{\widetilde{\chi}}}
\newcommand{\wttau}{{\widetilde{\tau}}}
\newcommand{\wtvp}{{\widetilde{\varphi}}}
\newcommand{\wtgamma}{{\widetilde{\gamma}}}
\newcommand{\wtGamma}{\widetilde{\Gamma}}
\newcommand{\wtlambda}{{\widetilde{\lambda}}}
\newcommand{\wtLambda}{{\widetilde{\Lambda}}}
\newcommand{\wtDelta}{{\widetilde{\Delta}}}
\newcommand{\wtOmega}{{\widetilde{\Omega}}}
\newcommand{\wtPsi}{{\widetilde{\Psi}}}
\newcommand{\wtPhi}{{\widetilde{\Phi}}}
\newcommand{\wtsigma}{{\widetilde{\sigma}}}
\newcommand{\wtxi}{{\widetilde{\xi}}}
\newcommand{\wtXi}{{\widetilde{\Xi}}}
\newcommand{\wtcA}{{\widetilde{\mathcal A}}}
\newcommand{\wtcB}{{\widetilde{\mathcal B}}}
\newcommand{\wtcC}{{\widetilde{\mathcal C}}}
\newcommand{\wtcD}{{\widetilde{\mathcal D}}}
\newcommand{\wtcK}{{\widetilde{\mathcal K}}}
\newcommand{\wtcH}{{\widetilde{\mathcal H}}}
\newcommand{\wtcN}{{\widetilde{\mathcal N}}}
\newcommand{\wtcP}{{\widetilde{\mathcal P}}}
\newcommand{\wtcQ}{{\widetilde{\mathcal Q}}}
\newcommand{\wtcR}{{\widetilde{\mathcal R}}}
\newcommand{\wtcS}{{\widetilde{\mathcal S}}}
\newcommand{\wtcT}{{\widetilde{\mathcal T}}}
\newcommand{\wtcU}{{\widetilde{\mathcal U}}}
\newcommand{\wU}{{\widetilde{\mathcal U}}}
\newcommand{\wtcV}{{\widetilde{\mathcal V}}}
\newcommand{\wtcZ}{{\widetilde{\mathcal Z}}}
\newcommand{\wtfN}{\widetilde{\mathfrak{N}}}
\newcommand{\wtfT}{\widetilde{\mathfrak{T}}}
\newcommand{\wtfU}{\widetilde{\mathfrak{U}}}
\newcommand{\wtfS}{\widetilde{\mathfrak{S}}}
\newcommand{\wtcM}{\widetilde{\mathcal M}}
\def\U{{\mathcal U}}
\newcommand{\fro}{{\mathfrak{r}}}
\newcommand{\dd}{{\mathfrak{d}}}
% widehat lowercase letters
\newcommand{\whb}{{\widehat b}}
\newcommand{\whd}{{\widehat d}}
\newcommand{\whu}{{\widehat u}}

% widehat uppercase letters
\newcommand{\whF}{{\widehat F}}
\newcommand{\whG}{{\widehat G}}
\newcommand{\whK}{{\widehat K}}
\newcommand{\whL}{{\widehat L}}
\newcommand{\whN}{{\widehat N}}
\newcommand{\whM}{{\widehat M}}
\newcommand{\whP}{{\widehat P}}
\newcommand{\whU}{{\widehat U}}
\newcommand{\whV}{{\widehat V}}
\newcommand{\whcD}{{\widehat{\mathcal D}}}
\newcommand{\whcH}{{\widehat{\mathcal H}}}
\newcommand{\whcR}{{\widehat{\mathcal R}}}
\newcommand{\whcN}{{\widehat{\mathcal N}}}
\newcommand{\whcP}{{\widehat{\mathcal P}}}
\newcommand{\whcM}{{\widehat{\mathcal M}}}
\newcommand{\whcT}{{\widehat{\mathcal T}}}
\newcommand{\whcX}{{\widehat{\mathcal X}}}
\newcommand{\whcZ}{{\widehat{\mathcal Z}}}
\newcommand{\whXi}{{\widehat{\Xi}}}
\newcommand{\wtwhXi}{\widetilde{\widehat{\Xi}}}
\newcommand{\whOmega}{{\widehat \Omega}}


% overline letters
\newcommand{\oA}{{\overline{A}}}
\newcommand{\oU}{{\overline{U}}}
\newcommand{\oV}{{\overline{V}}}
\newcommand{\Om}{\Omega}

% the rest of these are shortcuts for standard text formats

% blackboard bolds
%\newcommand{\B}{{\mathbb B}}
%\newcommand{\C}{{\mathbb C}}
\newcommand{\D}{{\mathbb D}}
\newcommand{\E}{{\mathbb E}}
\newcommand{\F}{{\mathbb F}}
%\newcommand{\G}{{\mathbb G}}
\newcommand{\K}{{\mathbb K}}
\newcommand{\N}{{\mathbb N}}
\newcommand{\Q}{{\mathbb Q}}
%\newcommand{\R}{{\mathbb R}}
\renewcommand{\S}{{\mathbb S}}
%\newcommand{\T}{{\mathbb T}}
%\newcommand{\Z}{{\mathbb Z}}

% cal script letters
\newcommand{\cA}{{\mathcal A}}
\newcommand{\cB}{{\mathcal B}}
\newcommand{\cC}{{\mathcal C}}
\newcommand{\cD}{{\mathcal D}}
\newcommand{\cE}{{\mathcal E}}
\newcommand{\cF}{{\mathcal F}}
\newcommand{\cG}{{\mathcal G}}
\newcommand{\cH}{{\mathcal H}}
\newcommand{\cI}{{\mathcal I}}
\newcommand{\cJ}{{\mathcal J}}
\newcommand{\cK}{{\mathcal K}}
\newcommand{\cL}{{\mathcal L}}
\newcommand{\cM}{{\mathcal M}}
\newcommand{\cN}{{\mathcal N}}
\newcommand{\cO}{{\mathcal O}}
\newcommand{\cP}{{\mathcal P}}
\newcommand{\cQ}{{\mathcal Q}}
\newcommand{\cR}{{\mathcal R}}
\newcommand{\cS}{{\mathcal S}}
\newcommand{\cT}{{\mathcal T}}
\newcommand{\cU}{{\mathcal U}}
\newcommand{\cV}{{\mathcal V}}
\newcommand{\cW}{{\mathcal W}}
\newcommand{\cX}{{\mathcal X}}
\newcommand{\cY}{{\mathcal Y}}
\newcommand{\cZ}{{\mathcal Z}}


% fraktur letters

\newcommand{\fA}{{\mathfrak{A}}}
\newcommand{\fa}{{\mathfrak{a}}}
\newcommand{\fB}{{\mathfrak{B}}}
\newcommand{\fb}{{\mathfrak{b}}}
\newcommand{\fC}{{\mathfrak{C}}}
\newcommand{\fD}{{\mathfrak{D}}}
\newcommand{\fg}{\mathfrak{g}}
\newcommand{\fG}{\mathfrak{G}}
\newcommand{\fH}{\mathfrak{H}}
\newcommand{\fh}{{\mathfrak{h}}}
\newcommand{\fk}{{\mathfrak{k}}}
\newcommand{\fK}{{\mathfrak{K}}}
\newcommand{\fl}{{\mathfrak{l}}}
\newcommand{\fm}{{\mathfrak{m}}}
\newcommand{\fL}{{\mathfrak{L}}}
\newcommand{\fM}{{\mathfrak{M}}}
\newcommand{\fN}{{\mathfrak{N}}}
\newcommand{\fO}{{\mathfrak{O}}}
\newcommand{\fP}{{\text{\gothfamily P}}}
\newcommand{\fp}{{\mathfrak{p}}}
\newcommand{\fQ}{{\text{\gothfamily Q}}}
\newcommand{\fq}{{\mathfrak{q}}}
\newcommand{\fR}{{\text{\gothfamily R}}}
\newcommand{\fr}{{\mathfrak{r}}}
\newcommand{\fS}{{\mathfrak{S}}}
\newcommand{\fT}{{\mathfrak{T}}}
\newcommand{\fU}{{\mathfrak{U}}}
\newcommand{\fV}{{\mathfrak{V}}}
\newcommand{\fX}{{\mathfrak{X}}}
\newcommand{\fY}{{\mathfrak{Y}}}
\newcommand{\fZ}{{\mathfrak{Z}}}

%text letters
\newcommand{\ta}{{\mathtt{a}}}
\newcommand{\tb}{{\mathtt{b}}}
\newcommand{\tc}{{\mathtt{b}}}
\newcommand{\td}{{\mathtt{d}}}
\newcommand{\te}{{\mathtt{e}}}
\newcommand{\tf}{{\mathtt{f}}}
\newcommand{\tg}{{\mathtt{g}}}

\newcommand{\ti}{{\mathtt{i}}}
\newcommand{\tj}{{\mathtt{j}}}
\newcommand{\tk}{{\mathtt{d}}}
%\newcommand{\tl}{{\mathtt{l}}}
\newcommand{\tm}{{(\mathtt{m})}}
\newcommand{\tn}{{\mathtt{n}}}

\newcommand{\tA}{{\mathtt{A}}}
\newcommand{\tB}{{\mathtt{B}}}
\newcommand{\tC}{{\mathtt{C}}}
\newcommand{\tD}{{\mathtt{D}}}
\newcommand{\tE}{{\mathtt{E}}}
\newcommand{\tF}{{\mathtt{F}}}
\newcommand{\tG}{{\mathtt{G}}}
\newcommand{\tH}{{\mathtt{H}}}
\newcommand{\tK}{{\mathtt{K}}}
\newcommand{\tI}{{\mathtt{I}}}
\newcommand{\tL}{{\mathtt{L}}}
\newcommand{\tM}{{(\mathtt{M})}}
\newcommand{\tN}{{\mathtt{N}}}
\newcommand{\tO}{{\mathtt{O}}}
\newcommand{\Tan}{{\cS_0}}


%bold letters
\newcommand{\ba}{{\bf a}}
\newcommand{\bb}{{\bf b}}
\newcommand{\bc}{{\bf c}}
%\newcommand{\bd}{{\vec \bf 2}}
\newcommand{\be}{{\bf e}}
\newcommand{\bi}{{\bf i}}
\newcommand{\bj}{{\bf j}}
\newcommand{\bk}{{\bf k}}
%\newcommand{\bm}{{\bf m}}
\newcommand{\bn}{{\bf n}}
\newcommand{\bu}{{\bf u}}
\newcommand{\bw}{{\bf w}}
\newcommand{\bv}{{\bf v}}
\newcommand{\bz}{{\bf z}}
\newcommand{\bN}{{\bf N}}
\newcommand{\bM}{{\bf M}}
%\usepackage{eufrak}
\usepackage{bm}


%wide check calligraphic letters
\newcommand{\wccT}{{\widecheck{\cT}}}


%greek letters
\newcommand{\oo}{{\omega}}
\newcommand{\al}{{\alpha}}
\newcommand{\bt}{{\beta}}


% differential operator
\renewcommand{\d}{\partial}
\newcommand{\dt}{\partial_t}
\newcommand{\lap}{\Delta}
\newcommand{\grad}{\nabla}
\newcommand{\derx}{\partial_x}
\newcommand{\derxx}{\partial_{xx}^2}
\newcommand{\suma}{\sum_{\substack{\al,\bt\in \N^\Z\\ |\al|=|\bt|\\ \sum_j j(\al_j-\bt_j)=0}}}
\newcommand{\w}[1]{\mathtt{w}_{#1}}
\newcommand{\ww}[1]{\mathtt{w}_{#1}\times \mathtt{w}_{#1}}

% norms
%\newcommand{\norm}[1]{\| #1 \|}
%\newcommand{\abs}[1]{\left| #1 \right|}
\newcommand{\0}{{(0)}}
%\newcommand{\2}{{(2)}}
%\newcommand{\1}{{(1)}}
\newcommand{\3}{{(3)}}
% braces
\newcommand{\la}{\left\langle}
\newcommand{\ra}{\right\rangle}

% mix
%\newcommand{\id}{\mathbbm{1}}
\newcommand{\im}{{\rm i}}
\newcommand{\jap}[1]{\langle #1 \rangle}
\newcommand{\re}[1] {\mbox{\Fontlukas T}\!_{#1}}  % {{\mathfrak T}\!_{#1}}
\newcommand{\und}[1]{\underline{#1}}
%\newcommand{\be}{{\bf e}}
\newcommand{\di}{{\rm d}}
\newcommand{\e}{{\varepsilon}}
\newcommand{\hh}{ h }


\newcommand{\pix}{{\pi_x}}
\newcommand{\piy}{{\pi_y}}
%\newcommand{\gen}{\fg}
\newcommand{\inv}{\partial_x^{-1}}
\newcommand{\uno}{{\mathbb I}}
\newcommand{\rdelta}{{\textcolor{red}{\delta}}}
\newcommand{\meas}{{\rm meas}}
\newcommand{\mix}{{\rm mix}}
\newcommand{\hor}{{\rm hor}}
\newcommand{\out}{{\rm out}}
\newcommand{\diag}{{\rm diag}}
\renewcommand{\line}{{\rm line}}
\newcommand{\fin}{{\rm fin}}
\newcommand{\Fjs}{{F^{(j,\s)}}}
\newcommand{\Fjsab}{{F^{(j,\s)}_{\al,\bt}}}
%\newcommand{\bu}{{\vec{\bf 1}}}
\newcommand{\bd}{{\vec{\bf 2}}}
%\newcommand{\red}[1]{{\textcolor{red}{#1}}}
\newcommand{\kam}{{\eta}}
\newcommand{\bnorm}[1]{{|\mkern-6mu |\mkern-6mu | \,  #1 \,  |\mkern-6mu |\mkern-6mu |}  }

\newcommand{\sQ}{{\mathscr{Q}}}
\newcommand{\sH}{{\mathscr{H}}}
\newcommand{\sD}{{\mathscr{D}}}
\newcommand{\sR}{{\mathscr{R}}}
\newcommand{\sS}{{\mathscr{S}}}
\newcommand{\tw}{{\mathtt{w}}}
\newcommand{\nnorm}[1]{{\left\vert\kern-0.25ex\left\vert\kern-0.25ex\left\vert #1 
    \right\vert\kern-0.25ex\right\vert\kern-0.25ex\right\vert}}
\newcommand{\coef}[1]{#1_{m,\alpha,\beta}}    %%%%%%prima c'era il b invece dell m
\newcommand{\coeffu}[1]{#1_{b^{1},\alpha^{1},\beta^{1}}}
\newcommand{\coeffd}[1]{#1_{b^{2},\alpha^{2},\beta^{2}}}
\newcommand{\bcoef}[1]{#1_{\alpha,\beta}}
\newcommand{\bcoeffu}[1]{#1_{\alpha^{1},\beta^{1}}}
\newcommand{\bcoeffd}[1]{#1_{\alpha^{2},\beta^{2}}}
\newcommand{\buno}{b^{1}}
\newcommand{\bdue}{b^{2}}
\newcommand{\aluno}{\alpha^{1}}
\newcommand{\aldue}{\alpha^{2}}
\newcommand{\btuno}{\beta^{1}}
\newcommand{\btdue}{\beta^{2}}
\newcommand{\uu}{{u^\alpha}{\bar{u}^\beta}}
\newcommand{\uuu}{u^{\aluno}\bar{u}^{\btuno}}
\newcommand{\uud}{u^{\aldue}\bar{u}^{\btdue}}
\newcommand{\buu}{{u^\bal}{\bar{u}^\bbt}}
\newcommand{\buuu}{u^{\baluno}\bar{u}^{\bbtuno}}
\newcommand{\buud}{u^{\baldue}\bar{u}^{\bbtdue}}
\newcommand{\pon}{{{\Pi^{0,\mathcal{K}}}}} %%%%% prima era \cN ma ora con le nuove notazioni e la storia del Taylor è proprio K
\newcommand{\por}{{\Pi^{0,\mathcal{R}}}}
\newcommand{\pd}{{\Pi^{-2}}}
\newcommand{\ono}[1]{#1^{0,\mathcal{N}}}
\newcommand{\oro}[1]{#1^{0,\mathcal{R}}}
\newcommand{\odo}[1]{#1^{-2}}
\newcommand{\modi}[1]{\abs{u_{#1}}^2}
\newcommand{\modu}{\abs{u}^2}
\newcommand{\La}[1]{\Lambda^{\pa{#1}}}
\newcommand{\es}{e^{\set{S,\cdot}}}
\newcommand{\bal}{{\bm \al}}
\newcommand{\bbt}{{\bm \bt}}
\newcommand{\baluno}{\bal^{1}}
\newcommand{\baldue}{\bal^{2}}
\newcommand{\bbtuno}{\bbt^{1}}
\newcommand{\bbtdue}{\bbt^{2}}
\newcommand{\Heta}{\cH_{r,s, \eta}}
\newcommand{\peseta}{e^{\eta\abs{\pi(\al - \bt)}}} %%%%% peso del momento 
\newcommand{\nore}[1]{\abs{#1}_{r,s,\eta}}
\newcommand{\noreq}[1]{\abs{#1}_{r,s,\eta}^{(q,a,p,\dep)} } %%% norma s,r,eta
\newcommand{\Lnor}[1]{\abs{#1}^{\rm Lip}}  %%% seminorma lipschitz
\newcommand{\gnor}[1]{\abs{#1}^{\gamma}}  %%% norma gamma
\newcommand{\Dom}[1]{\sum_{#1} \omega_{#1} \pa{\abs{u_{#1}}^2 - I_{#1}}}  %% D_omega con l'indice che voglio
\newcommand{\azioj}[1]{\pa{\abs{u_{#1}}^2 - I_{#1}}} %%% coef azioni con indici 
\newcommand{\azio}{\pa{\abs{u}^2 - I}} %% coeff azione
\newcommand{\divisor}[1]{#1\cdot\pa{\bal - \bbt}}
\newcommand{\zero}[1]{#1^{(0)}}
\newcommand{\zeroR}[1]{#1^{(0,\cR)}}
\newcommand{\zeroK}[1]{#1^{(0,\cK)}}
\newcommand{\due}[1]{#1^{(-2)}}
\newcommand{\buon}[1]{#1^{\ge 2}}
\newcommand{\costa}{C e^{C(\teta,p){\pa{\frac{\sigma}{2}}}^{-3/\teta}}}  %% la costante del lemma eqz homologica
\newcommand{\call}{\mathcal{L}}
\newcommand{\esp}[2]{e^{\set{#1, #2}}}
\newcommand{\red}[1]{\textcolor{red}{#1}}
\newcommand{\hv}{\hat{v}}
\newcommand{\hp}{\hat{\varphi}}
\newcommand{\hu}{u^\eta}

\newcommand{\blue}[1]{\textcolor{blue}{#1}}
%%%%%%%%%%%%%%%%%%%%%% fine macros %%%%%%


\RequirePackage{url}
\newtheorem{prop}{Proposition}[section]
    \newtheorem{thm}{Theorem}[section]
    \newtheorem{ex}{Example}
    \newtheorem{cor}{Corollary}
\newtheorem{rmk}{Remark}[section]
\newtheorem{defn}{Definition}
\newtheorem{lemma}{Lemma}
\newtheorem{corollary}{Corollary}
\newtheorem*{Theorem*}{Theorem}

%%%%%% numerazione %%%%%%%
\numberwithin{equation}{section}
%\numberwithin{thm}{section}
\numberwithin{defn}{section}
\numberwithin{prop}{section}
\numberwithin{cor}{section}
\numberwithin{lemma}{section}
\numberwithin{rmk}{section}
\numberwithin{corollary}{section}
\newcommand{\To}{\operatorname{T}_0}
\usepackage{showkeys}
%%%%%%%%%%
%
\begin{document}
\date{\today}
\author{Bassam Fayad and Jessica Elisa Massetti}


\title[symplectic Vs dissipation]{Normal forms for perturbations of general dissipative vector fields stemming from Hamiltonian ones ??}
\begin{abstract}

\end{abstract}

\maketitle
\tableofcontents

\chapter{ciao}
\section{Introduction and main results}


Let $\cH$ be the space of germs of real analytic Hamiltonians defined in a neighborhood of $\To :=\T^n\times\set{0}$ and let $\Vh$ be the corresponding space of Hamiltonian vector fields.

- For $\alpha\in\R^n$, let $\cK^\alpha$ be the affine subspace of $\cH$ consisting of \emph{non degenerate} Hamiltonians defined in a neighborhood of $\To$, $$\mathcal{K}^\alpha=\set{K\in\mathcal{H}: K(\teta,r)= c + \alpha\cdot r + \frac{1}{2}Q(\teta)\cdot r^2 + O(r^3)},$$

where $Q$ is a non-degenerate quadratic form defined on $\T^n$, which we shall refer to as \emph{twist}.
For any $K\in\cK^\alpha$, let us denote by $X_K$ the associated Hamiltonian flow 
\begin{equation}\label{Ham vf}
	X_K=\eqsys{\dot\teta = \frac{\partial K}{\partial r}(\teta,r)= \alpha + O(r)\\
		\dot{r} = -\frac{\partial K}{\partial \teta}(\teta,r)= O(r^2).} 
\end{equation}
and denote by $\Uh(\alpha)$ the corresponding affine subspace of $\Vh$ of Hamiltonian vector fields of the form \eqref{Ham vf}.

Note that a Hamiltonian in $\mathcal{K}^\alpha$  has the torus $\normal{T}^n_0$  invariant by the associated Hamiltonian flow, which restricted to that torus is the translation by  frequency $\alpha$, $\teta \mapsto \teta + \alpha$.
In the following we shall assume that the twist $Q$, satisfies $\left|\det\frac{1}{(2\pi)^n}\int Q(\teta)\,d\teta\right| \geq 1/100$.
%-quasi-periodic and has a twist condition around the torus. 
%We want to explore the persistence of a QP torus with frequency $\alpha$ under addition of dissipation terms and perturbation.

\medskip

- Let now $\eta \in [-1,1]$  and and let us consider the extension $\cW$ of $\Vh$,  cosisting of vector fields $$w = \vh + v^\eta = \vh + \eta v$$ obtained just by adding non-necessarily-Hamiltonian directions that exhibits a factor $\eta$ as $\eta v$, which we denote by $\cV^\eta$.  
\begin{equation}\label{hybrid vf}
\cW := \Vh\times\cV^\eta .
\end{equation}

Note that the non necessarily hamiltonian vector field $v$ may depend (smoothly) on $\eta$ as well.

- In line with the above definition,  given $\eta\in [-1,1]$ and $M \in {\rm GL}(n,\R)$ diagonalizable with real eigenvalues $\mu_j$'s, let $\cU^\eta(M)$ be the affine subspace of $\cV^\eta$ of real analytic vector fields on $\T^n\times\R^n$ of the form  
\begin{equation}
	u^\eta = \eqsys{\dot{\teta} =  \eta O(r)\\
		\dot r = \eta M\cdot r +\eta O(r^2) \,. }
	\label{u}
\end{equation}

so that, given $\alpha\in\R^n$ 
\begin{equation}\label{good U}
	\U(\alpha,M,\eta) := \Uh(\alpha)\times\cU^\eta(M)
\end{equation}
is the affine subspace of good kind of vector-fields, i.e. those  possessing the invariant hyperbolic quasi-periodic torus $\To$.

- \textit{$n$ parameters families of v.f.}: 
Let now $\nu \in \R^n$ be a vector of free parameters, we may consider the following $n$-parameter family of vector spaces  {$$\cW(\nu) = \cW\oplus (\eta\nu\partial_r)\,,\quad \nu\in\R^n$$} consisting of vector fields stemming from $\cW$, of the form 
\begin{equation}
	w + \eta \nu\partial_r\, \quad \quad \nu\in\R^n
	\end{equation} 



\begin{rmk}[v.f. with translated quasi-periodic hyperbolic torus]
 The (family of) affine subspaces of $\cW(\nu)$  parameterized by $\nu\in\R^n$ 
\begin{equation}
	\label{quasi-good U}
\cU(\alpha,M,\eta;\nu) :=\U(\alpha,M,\eta)\oplus(\eta\R^n\partial_r)
\end{equation}
consists of those vector fields of the form
 $$
 U = X_K + u^\eta + \eta\nu\partial_r\,\quad \quad \nu\in\R^n\,,\quad K\in\cK^\alpha\,,\quad u^\eta\in\cU^\eta(M)
 $$
 i.e. those ones whose flow translates $T_0$ in the normal directions by $\eta\nu$.
 \end{rmk}


\begin{rmk}
Both parameters $\eta,\nu$ twist the Hamiltonian structure of $\Uh(\alpha)$. In particular, note that, when $\nu = 0$,  for a vector field  $U\in \U(\alpha,M,\eta; 0)$, the torus $\normal{T}^n_0 = \T^n\times\set{0}$ is invariant, $\alpha$-quasiperiodic, and hyperbolic with normal dynamics given by the Floquet exponents $\eta \mu_1,\ldots,\eta \mu_n,$ the eigenvalues of $\eta M$. Otherwise $\normal{T}^n_0$ is translated by $\eta\nu$ in the normal direction. \\

On the other hand, when both $\eta=0$ and $\nu=0$, the space $\cU(\alpha,M, 0;0)\equiv\cU(\alpha,0,0;0) \equiv \Uh(\alpha)$ coincides with the purely Hamiltonian one.
\end{rmk}


- Let $\Lambda $ be the subspace of $\cW$ of vector fields of the form
\begin{equation}
	\lambda(\teta,r) = (0,  b +  B\cdot r)\,,\quad b\in\R^n\,, B\in\Mat_n(\R^n)
	\end{equation}
we shall denote with $\lambda^\eta$ elements in $\Lambda$ of the form 
$$
\lambda^\eta(\teta,r) = (0,  \eta b + \eta B\cdot r)
$$ and the corresponding subset $\Lambda^\eta$.  \\

\red{BASSAM: ajouter un commentaire sur le fait que, comme on veut une transformation du type $\id + O(\eta)$ t.q. tout devien smooth hamiltonien quand $\eta\to 0$, il n'y a pas la condition $[B,M] = 0$ mais il faut une matrice $B$ pleine}\\ 
\begin{itemize}[leftmargin=*]
\item  Let $\cG$ be the space of germs of real analytic diffeomorphism of $\T^n\times\R^n$ of the form 

\begin{equation}\label{diffeo dissip}
	g(\teta,r) = (\varphi(\teta), R_0(\teta) + R_1(\teta)\cdot r)\,,
\end{equation} 
where $\varphi = \id + v\,, v:\T^n\to\R^n$ is a diffeomorphism of the torus fixing the origin, while $R_0, R_1$ are  $\R^n$ and $\Mat_n(\R^n)$-valued functions respectively such that $R_0(0)=0$ and $R_1(0) - I = 0$.\\
As before, we shall denote by $\cG^\eta$, the subspace of diffeomorphisms in $\cG$ of the form 
$$
g^\eta (\teta,r) = (\varphi^\eta(\teta), R^\eta_0(\teta) + R^\eta_1(\teta)\cdot r) = (\teta + \eta v(\teta), r + \eta R_0(\teta) +\eta (R_1(\teta) - I)\cdot r)
$$

\item Let $\cG^\omega$ be the space of germs of real analytic symplectic diffeomorphisms of $\T^n\times\R^n$ of the form 
\begin{equation}\label{diffeo symp}
	h(\teta,r) = (\varphi(\teta), (r + dS(\teta) + \xi)\cdot\varphi'^{-1}(\teta))
\end{equation}
where $\varphi$ is a diffeomorphism as above and $S:\T^n\to \R\,, S(0)=0$  \\
 and define the product space 
$$
\hat \cG := \cG^\omega\times\cG \, \, \supset\,\,  \cG^\omega\times\cG^\eta  .
$$

Note that as $\eta\to 0$, the space $\Lambda^\eta$ reduces to the null space while $\cG^\omega\times\cG^\eta$ to the one of symplectic diffeomorphisms.
\end{itemize}

- In order to handle small divisors, we suppose non-resonant condition for the tangent frequencies. The spectrum of $M$ being real, it will suffice to suppose the following classic arithmetic condition, onvolving only the frequency of the linear flow on the torus.\\
 For $\alpha\in\R^n$,%, $\mu \in \R^n$, 
	$\tau,\gamma>0$, 
	we say that $\alpha \in DC(\tau,\gamma)$ if the following Diophantine condition holds
	\begin{align}
		\abs{k\cdot\alpha}&\geq \frac{\gamma}{\abs{k}^\tau}\quad\forall k\in\Z^n\setminus\set{0},\end{align}
	where we have set $\abs{k} = \abs{k_1} + \ldots + \abs{k_n}$.\\

In general, for the Floquet exponents $\mu_j$'s, we shall only assume that they are pairwise distinct and different from $0$, with a uniform fixed lower bound 

\begin{equation}\label{mu conditions}
    \min_{i\neq j} |\mu_i - \mu_j| > \gamma\,, \quad \quad |\mu_j| > \gamma \quad \forall i=1,\ldots n\,.
\end{equation}


%$\bullet$ Observe that for $U = X_K + V_\nu$ the normal form is tautological with $\Phi = \id$, $U = X_K + V_0$ and $\nu' = \nu$.
%
%\medskip
%
%\begin{rmk}[previous]
%$\bullet$ Note that by symplectic change of variable that shifts $r$ (and conjugate $Q$ to its average part), we can for any $\beta \in B(\alpha,\eps_0)$ (where $\eps_0\ll |\det(\widehat{Q})|$), find a conjugacy $\Psi_\beta$  such that  $${\Psi_{\beta}}_*W \in \U(\beta,M', \eta;\nu')$$
%with $M',Q'$ are $\sqrt{\eps_0}$-close to $M, Q$, and $\nu'$ is $\sqrt{\eps_0}$-close to  $\nu+M\widehat{Q}^{-1}(\beta-\alpha)$. 

%Since $\nu'(\beta)\sim \nu+M\widehat{Q}^{-1}(\beta-\alpha)$. Hence there is a function $\beta(\nu)$ such that $$\left(X_{K+\eps H}+u^\eta+\epsilon v^\eta\right)\circ 
	%\Phi \in \U(\beta(\nu),M',\eta;0).$$
%Moreover, $\partial_\nu \beta \sim \widehat{Q}M^{-1}$. 
%Hence, for a large measure set of $\nu$, $\beta(\nu)\in 
%{\rm DC}(\tau,\gamma)$.
%\end{rmk}

%\red{Why $\sqrt{\eps_0}$ ?}
As is customary, in Section \ref{complex extensions} we will introduce complex extensions of a certain width $s>0$ of manifolds and define the corresponding
spaces of the real analytic vector fields and diffeomorphisms introduced above, enjoying such an extension. We will endow such spaces
with a Banach norm. Thus, all the closeness conditions appearing in the statements of this
section have to be understood as referring to that norm (see formula \eqref{W norm}).  

\medskip

%%%%%%%%%%%%%%%%%%%%%%%%%  teo di prima %%%%%%%%%%%%%%%%%%%%%%%%%%%%%%

%\begin{thm}[Normal form 2]\label{th.NF2} Fix  $\tau,\gamma>0$. There exists $\eps>0$ such that the following holds. Let $\a \in {\rm DC}(\tau,\gamma)$. Assume that $U\in \U(\alpha,Q,\eta,M,\nu)$ with $\eta \in [-1,1]$, $Q$ a non degenerate quadratic form $Q$ : $\left|\det\frac{1}{(2\pi)^n}\int Q(\teta)\,d\teta\right| \geq 1/100$,   $M \in {\rm GL(n,\R)}$ diagonalizable on $\R$ with $\min_{i\neq j}|\mu_i-\mu_j|>1/100$, $\min |\mu_i|>1/100$.  For any $\beta \in {\rm DC}(\tau,\gamma) \cap B(\alpha,\zeta)$, there exist $M',Q',\nu'$  and a  unique affine conjugacy $\Phi=\Id+\phi_{\text{Ham}}+\eta \psi$
%%, and a  vector $\Lambda=\eta(0,\lambda)\in \T^d \times \R^d$,
%such that $$\left(X_{K+\eps h}+V+\epsilon \eta v\right)\circ 
%	\Phi \in \U(\beta,\eta,M',Q',\nu')$$
%	with $M'$ and $Q'$  $\sqrt{\eps_0}$-close to $M$ and $Q$ and $\nu'$ is $\sqrt{\eps_0}$-close to  $\nu+M\widehat{Q}^{-1}(\beta-\alpha)$
%\end{thm}

%%%%%%%%%%%%%%%%%%%%%%%%%%%%%%%%%%%%%%%%%%




\medskip


%%%%%%%%%%%%%%%%%% teo di prima %%%%%%%%%%%%%%%%%%%

%\begin{thm}[Normal form with counter terms]
%Fix  $\tau,\gamma>0$. There exists $\eps>0$ such that the following holds. Let $\a \in {\rm DC}(\tau,\gamma)$. Assume that $U\in \U(\alpha,M, \eta;\nu)$ with $\eta \in [-1,1]$,  $M \in {\rm GL(n,\R)}$ diagonalizable on $\R$ with $\min_{i\neq j}|\mu_i-\mu_j|>1/100$, $\min |\mu_i|>1/100$.
%\smallskip
%  For any $\al' \in {\rm DC}(\tau,\gamma) \cap B(\alpha,\zeta)$ and any $M'\in\GL(n,\R)$ diagonalizable with $\min_{i\neq j}|\mu'_i-\mu'_j|>1/100$, $\min |\mu_i|>1/100$ , there exist $B,Q'$ and $\beta, b\in\R^n$  and a  unique affine conjugacy $\Phi=\Id+\phi_{\text{Ham}}+\eta \psi$
%%, and a  vector $\Lambda=\eta(0,\lambda)\in \T^d \times \R^d$,
%such that $$\left(X_{K+\eps h}+V+\epsilon \eta v\right) = \Phi_{*} U + (0, B\cdot r + b) \quad \quad U\in \U(\a',\eta,M',Q',0)$$
%	with $B$ a diagonal matrix, $Q'$  $\sqrt{\eps_0}$-close to $M, \a$ and $Q$ respectively, while $b$ is $\sqrt{\eps_0}$-close to  $\nu+M\widehat{Q}^{-1}(\al'-\alpha)$
%\end{thm}

%%%%%%%%%%%%%%%%%%%%%%%%%%%%%%%%%%%%%%%


The cornerstone of the present work is the following normal form theorem. 

\begin{thm}[\blue{Main normal form theorem}]
	\label{Main NFT}
	 Let $\tau,\gamma > 0$, let $\alpha\in DC(\tau,\gamma)$ and $M\in\Mat_n(\R)$ be diagonalizable with pairwise distinct eigenvalues $\mu_j$'s satisfying condition \eqref{mu conditions}, and let $u^0 = (X_{K^0},u^{0,\eta})\in\cU(\alpha,M,\eta)$. If $w\in\cW$ is sufficiently close to $u^0$, then there exists a unique triplet $(g,u,\lambda)\in\hat\cG\times\cU(\alpha,M,\eta)\times \Lambda$, where $g=(h,g^\eta)\in\cG^\omega\times\cG^\eta$, is close to $((\id,\id), u^0,0)$, such that $g=(h,g^\eta),u=(X_K, u^\eta), $ and $ \lambda^\eta$ satisfy
	\begin{equation}
w = h_* g^\eta_* u  + \lambda^\eta\,.
	\end{equation}
	
\end{thm}


\begin{rmk}
Observe that given $$U = (\alpha + Q\cdot r, \eta Mr + \eta\nu),$$ by adding the counter term   $$\lambda= b + B\cdot r = -\eta\nu + \eta M'Q^{-1}(\alpha'-\alpha)+ \eta(M' - M)\cdot r$$ and performing the (symplectic) change of variables $T: r\mapsto r + Q^{-1}(\alpha - \alpha')$, we push $U$ forward to the normal form $$T_* U = (\alpha' + Q r, \eta M' r ), $$ where the torus $\To$ is $\alpha'$-quasi periodic, with Floquet exponents governed by the eigenvalues of $M'$.
\end{rmk}


The following parameter-dependent result will follow straightforwardly. 

\begin{thm}[Normal form with counter terms for $\nu$-families of v.f.]\label{thm NF parameters}
Fix  $\tau,\gamma>0$. Let $\a \in {\rm DC}(\tau,\gamma)$ and $M \in {\rm GL(n,\R)}$ be diagonalizable on $\R$ with pairwise distinct eigenvalues $\mu_j's$ satisfying \eqref{mu conditions}.
Assume that 
$$
u^0 = (X_{K^0}, u^{0,\eta}) \,\,\in \U(\alpha,M, \eta)
$$ with $\eta \in [-1,1]$. 
 Then, there exists $\eps>0$ such that the following holds.

\smallskip
  For any $\al' \in {\rm DC}(\tau,\gamma) \cap B(\alpha,\zeta)$ and any $M'\in\GL(n,\R)$ diagonalizable with pairwise distinct eigenvalues $\mu'_js$ satisfying \eqref{mu conditions}, if $w=(\vh, v^\eta + \eta\nu)\in\cW(\nu)$ is $\eps$-close to $u^0 = (X_{K^0}, u^{0,\eta}) $, then there exist a unique triplet $(g,u,\lambda^\eta)\in\cG\times\cU(\alpha',M',\eta)\times \Lambda^\eta$, where $g=(h,g^\eta)\in\cG^\omega\times\cG^\eta$, is close to $(\id,\id)$, $u$ close to $u^0$ and $\lambda^\eta = (0, \eta b + \eta B\cdot r)$  close to $(0,\eta\nu)$  such that 
  \begin{equation}
  	w = h_* g^\eta_* u  + \lambda^\eta\,,\quad \quad u = X_K + u^\eta\,.
  \end{equation}
In particular, $B$ is a diagonalizable matrix,  $\sqrt{\eps}$-close to $M'$, while $\eta b$ is $\sqrt{\eps}$-close to  $\eta\nu+\eta M'\widehat{Q}^{-1}(\al'-\alpha)$ \red{I am not sure about the $\sqrt{\eps}$ while not $O(\eps)$}
\end{thm}


\begin{rmk}
 \blue{Note that differently from Moser's normal form counter terms $b, B$ are chosen to kill all the average part of the r.h.s. in the homological equations, and not just the kernels terms of $M$ and $[M,\cdot]$. Therefore that we are not left with denominators like $\frac{1}{\eta(\mu_i - \mu_j)}$ and we can take smoothly the hamiltonian limit $\eta\to 0$ } 
 \end{rmk} 

Note that, in particular, $B = B(M'), b = b(\nu, \alpha')$ depend (Whitney) smoothly on the frequencies/parameters. 

\medskip

\subsection{Idea of the proof of Main Theorem \ref{Main NFT}: Kolmogorov and the Main bridge theorem}

The following formulation of the problem is in the spirit of an inverse function theorem à la Hamilton-Zehnder (see \cite{Zehnder:1975}) on scales of Banach spaces, with the notations and general setting proposed in \cite{Massetti:ETDS,Fejoz:2015}.

Because of the structure of the vector fields involved, provided the perturbation is sufficiently small, we shall apply Kolmogorov theorem as a black-box, in order to put our vectorfield in a preliminary normal form that already has the Hamiltonian part redressed. The remaining concern would then be to straighten the leftovers of order $O(\eta)$ by a dissipative, close to identity diffeomorphism $g\in\cG^\eta$.


%The conjugacy that realizes equation \eqref{final normal f}
%\begin{equation*}
% (K_1+V+\epsilon \eta v-\Lambda\circ \Phi_1)\circ \Psi \in \U(\alpha,\eta,M',Q')
% \end{equation*}
% is given by $\Phi_1\circ\Psi$, where $\Phi_1$ is the symplectic transformation given by Kolmogorov theorem.

\medskip

Let $w \in \cW$, i.e. a vectorfield of the form 
\begin{equation}\label{perturbed vf}
w= X_H + \eta v = \vh + v^\eta\,, \quad \textit{with $H$  close to $K^0$ and $v^\eta$ close to $u^{0,\eta}$.}
\end{equation}

be given.

\medskip

{Since $X_{K^0}$ is non-degenerate, and $\alpha$ Diophantine, if $\vh$ is sufficiently close to $X_{K^0}$, by Kolmogorov theorem, there exist a symplectic affine diffeomorphism $h\in\cG^\omega$ and $K$ a Hamiltonian in normal form,  such that 
\begin{equation}\label{hamiltonian conj}
\vh =  h_*X_K
\end{equation}
 with $K\in\cK^\alpha$, where $(K, h)$ is in some  neighborhood of $({K^0}, \id)\in\cK^\alpha\times\cG^\omega$. }  \\

More specifically,

\begin{thm}[Kolmogorov] 
\label{Kolmogorov}
The map
 $$
 \Phi_{\tK}: \Uh_{s+\sigma}(\alpha) \times \cG^{\omega}_{s+\sigma} \to \Vh_{s} \quad (X_{K},h)\mapsto h_{*}X_{K} 
 $$
 is a local diffeomorphism in the neighborhood of $(X_{K^{0}},\id)$ in the sense that 
 for any $\delta< s <s+\sigma<1$ there exist $\eps_{1}>0$ and a unique $C^\infty$-map $\Psi_\tK$ $$\Psi_\tK : \mathcal{V}^{\operatorname{Ham}, \eps_{1}}_{s+\sigma}\to \G^{\omega,\delta}_{s}\times\Uh_s(\alpha)$$ such that $\Phi_{\tK}\circ\Psi_\tK = \id\,,$
 where  $ \mathcal{V}^{\operatorname{Ham}, \eps_{1}}_{s+\sigma}$ stands for the ball of radius $\eps_1$ in $\Vh_{s+\sigma}$ centered at $X_{K^0}$ and $\G^{\omega,\delta}_{s}$ the one of radius $\delta$ centered at the identity in $\G^\omega_s$.
 \end{thm}
 

The equation above entails that $h(T^0)$ is invariant and $\alpha$-quasi periodic for $\vh$.  The solutions are real analytic on a domain of width $s'< s$. \\
 
 The domain of application of the theorem then provides us with the appropriate small $\eps_1$-ball of $X_{K^0}$ for $\vh$. 




The idea is to apply the above theorem to the vector field $w$ defined in \eqref{perturbed vf}, provided its norm is accordingly small, so that the Hamiltonian side of the perturbations is straightened to a normal form (that of course would depend on the conjugacy $h$ and the Hamiltonian terms $(O(r),O(r^2))$ of $X_K$. \\
The proof of Main theorem \ref{Main NFT}, will be articulated in two main steps

 
 \medskip
 
1) \textbf{Kolmogorov step:} In $\cW_{s+\sigma}$, let us so consider the open ball of radius $\eps_{1}$ centered at $w^{0} := X_{K^{0}} + u^{0,\eta}$, that is $B_{s+\sigma}^{\eps_{1}} := \set{w\in\cW_{s+\sigma}\,:\, \|w - w^0\|_{{s+\sigma}} < \eps_{1}}$. We shall say that a vector field $w\in\cW_{s+\sigma}$ is \emph{admissible} if $w\in B_{s+\sigma}^{\eps_{1}} $. 

% Let us now consider the map $\Psi_{1}: B_{s+\sigma}^{\eps_{1}} \to  $
Hence, to any admissible $w = \vh + \eta v$, by Kolmogorov theorem, there corresponds a couple $(X_{K}, h) = \Psi_\tK(\vh)$ such that $\vh = h_*X_K \,\,\Leftrightarrow\,\, h^{*}\vh = X_{K}$.

\smallskip
\smallskip

2) \textbf{Dissipative Step:} We now use $h$ to pull-back the vector field $w$ in \eqref{perturbed vf} and consider the $h$-dependent  $$h^*w= h^*\vh + h^*v^\eta = X_{K} +  h^*v^\eta,\quad \quad K = H\circ h$$

 {So the problem becomes, in a neighborhood of $h(\To)$ }, if  $v^\eta$  is close enough to $u^{0,\eta}$ (hence $h^{*}v^\eta$ is close to $h^*u^{0,\eta}$) we want to find \\
- a dissipative affine diffeomorphism  $g^\eta\in\cG^\eta$ (i.e. such that $g^\eta - \id = O(\eta)$)\\
 - $u^\eta\in\cU^\eta(M)$\\
 - $\lambda^\eta = (0, \eta b + \eta B\cdot r)$ \\
 such that 
\begin{equation}\label{dissipative normal f}
	h^*(\vh + v^\eta - \lambda^\eta ) = g_*(X_K + u^\eta) \, \quad \, \Leftrightarrow \, \quad \vh +  v^\eta = h_*g^\eta_*(X_K + u^\eta) + \lambda^\eta\,.
	\end{equation}
 \\


\medskip


Hence the torus $h\circ g^\eta (T^0)$ is invariant for $\vh + v^\eta - \lambda$ . Note that as $\eta\to 0$, we get the original Kolmogorov normal form \eqref{hamiltonian conj}. 

\medskip

Therefore for any given $h, X_K, \vh = \Phi_{\tK}(X_K, h)$, the problem then amounts to proving the following 

\begin{thm}[Bridge Normal form in a neighborhood of $h(T_0)$]\label{bridge thm}
    For any fixed $(h,K)$ in the domain of application of Kolmogorov Theorem \ref{Kolmogorov}, 
    the map
    \begin{equation}
\label{phi intermedio} 
\phi_{h,K}:   \cG^\eta\times\cU^\eta(M)  \times \Lambda^\eta \to h^{*}\cW\, \quad  ( g^\eta, u^\eta, \lambda^\eta)\mapsto  g^\eta_*(X_K + u^\eta) + h^*\lambda^\eta
\end{equation}
is locally invertible in a neighborhood of $(\id, u^{0,\eta},0)$, for any $h^* v^\eta$ sufficiently close to $h^*u^{0,\eta}$. Hence equation $$h^*(\vh + v^\eta) = g^\eta_*(X_K + u^\eta) + h^*\lambda^\eta$$ is solved in a neighborhood of $g(\To)$, via a unique $(g^\eta, u^{\eta}, \lambda^\eta)$, appropriately close to $(\id, u^{0,\eta}, 0)$.
\end{thm}

If $h^*v^\eta$ is small enough, and one achieve in eliminating $\lambda^\eta$, then $h\circ g^\eta(\To)$ is an invariant torus for $w$, with characteristic frequencies $(\alpha, \eta M)$.

\subsection{Elimination of parameters and invariant tori results}



\begin{thm}[Invariant tori (consequence of elimination of parameters)]\label{th.NF3} If $\eps>0$ is sufficiently small for a set of positive measure of $\nu$,  the vector field
$w$ $\eps$-close to $u^0$ has an invariant quasi-periodic torus, that is normally hyperbolic with Floquet exponents close to $\eta \mu_1,\ldots,\eta \mu_n,$ the eigenvalues of $\eta M$.
\end{thm}

\blue{
The above theorem will follow from the fact that the map 
$$ \left(\rm{DC}(\gamma,\tau)\times \GL_n(\R^n), ( \nu, M)\right) \, \longrightarrow \, \left(\R^n\times \GL_n(\R^n), (0,0)\right) 
, \quad (\a',M') \, \mapsto (b, B)
$$
is a local diffeomorphism.\\
$\bullet$ Elimination of $B$: Considering the map $\eta M\mapsto B(M)$, where $M$ has pairwise distinct real eigenvalues $\mu_{j}$ and $\partial_{\mu}B(M)$ is close to identity, one can determine $M'$ close to $M$ such that $B(M') =0$. This results in a translated torus theorem: $h\circ g(\To)$ is translated by $\eta b\in\R^{n}$ in the action direction. \\
$\bullet$ Eliminaton of $b = b(\alpha', \nu)$: by moving $\alpha'$ or $\nu$.}

\medskip


\begin{thm}[Invariant tori 2]\label{th.NF4} Consider an integrable twist Hamiltonian 
$$H(r,\theta)=H_0(r)$$
Add to its vector field $X_H$ a dissipation term 
$$X=X_H+\eta M (r-\nu)$$
$M$  diagonalizable on $\R$ with $\min_{i\neq j}|\mu_i-\mu_j|>1/100$, $\min |\mu_i|>1/100$. 

For $h(r,\theta)$ and $v(r,\theta)$ bounded by $1$ consider 
$$X_\eps:=X_{H+\eps h}+\eta M (r-\nu)+\eta \eps v$$
Then we have that for $\eps>0$ sufficiently small, for any $\eta \in [-1,1]$, for a positive measure set of $\nu \in \R^d$ that goes to full measure as $\eps \to 0$, the flow of $X_\eps$ has normally hyperbolic Diophantine quasi-periodic tori with Floquet exponents close to $\eta \mu_1,\ldots,\eta \mu_n,$ the eigenvalues of $\eta M$.
\end{thm}


\rule{\textwidth}{0.4pt}


\rule{\textwidth}{0.4pt}


\section{Functional setting}
\subsection{Complex extensions}
\label{complex extensions}
Let us extend the tori \[\T^n=\R^n/{2\pi\Z^n}\qquad \text{and}\qquad \normal{T}^n_0=\T^n\times\set{0}\subset\T^n\times\R^m,\] as \[\T^n_{\C}= \C^n/{2\pi\Z^n}\qquad \text{and} \qquad  \text{T}^n_\C = \T^n_{\C}\times\C^m\]  and, letting $s>0$, consider the corresponding $s$-neighborhoods defined using $\ell^\infty$-balls (in the real normal bundle of the torus): \[\T^n_s=\set{\teta\in\T^n_{\C}:\, \max_{1\leq j \leq n}\abs{\Im\teta_j}\leq s}\quad\text{and}\quad \text{T}^n_s=\set{\pa{\teta,r}\in\text{T}^n_{\C}:\, \abs{(\Im\teta,r)}\leq s},\]
where $\abs{(\Im\teta,r)}:= \max\pa{\max_{1\leq j\leq n}\abs{\Im\teta_j},\max_{1\leq j\leq m}\abs{r_j}}$. \\

Let now $f: \normal{T}^n_s\to \C$ be holomorphic, and consider its Fourier expansion $f(\teta,r)=\sum_{k\in\Z^n}\,f_k(r)\,e^{i\,k\cdot\teta}$, noting $k\cdot\teta = k_1\teta_1 +\ldots + k_n\teta_n$. In this context we introduce the so called "weighted norm": 
\begin{equation}\label{W norm}
	\abs{f}_s := \sum_{k\in\Z^n}\abs{f_k}\, e^{\abs{k}s},\quad \abs{k}=\abs{k_1}+\ldots +\abs{k_n},
\end{equation}
where $\abs{f_k}=\sup_{\abs{r}<s} \abs{f_k(r)}$. Whenever $f : \normal{T}^n_s\to\C^{n}$, $\abs{f}_s = \max_{1\leq j\leq n}(\abs{f_j}_s)$, $f_j$ being the $j$-th component of $f(\teta,r)$.\\ It is a trivial fact that the classical sup-norm is bounded from above by the weighted norm: \[\sup_{z\in{\normal{T}^n_s}}\abs{f(z)}\leq\abs{f}_s\] and that $\abs{f}_s<+\infty$ whenever $f$ is analytic on its domain, which necessarily contains some $\normal{T}^n_{s'}$ with $s'>s$. In addition, the following useful inequalities hold if $f,g$ are analytic on $\normal{T}^n_{s'}$ \[\abs{f}_s\leq\abs{f}_{s'}\,\text{ for }\, 0<s<s',\] and \[\abs{fg}_{s'}\leq \abs{f}_{s'}\abs{g}_{s'}.\] Moreover, one can show that if $f$ is analytic on $\normal{T}^n_{s+\sigma}$ and $g$ is a diffeomorphism of the form \eqref{diffeo dissip} sufficiently close to the identity, then $\abs{f\circ g}_s \leq C_g \abs{f}_{s+\sigma}$, where $C_g$ is a positive constant depending on $\abs{g - \id}_s$. For more details about the weighted norm, see for example \cite{Meyer:1975, Chierchia:2003}.\\
In general for complex extensions $U_s$ and $V_{s'}$ of $\T^n\times\R^m$, we will denote $\mathcal{A}(U_s,V_{s'})$ the set of real holomorphic functions from $U_s$ to $V_{s'}$ and $\mathcal{A}(U_s)$, endowed with the $s$-weighted norm, the Banach space $\mathcal{A}(U_s,\C)$.

Eventually, let $E$ and $F$ be two Banach spaces,

\begin{itemize}[leftmargin=*]
	\item We indicate contractions with a dot "$\,\cdot\,$", with the convention that if $l_1,\ldots, l_{k+p}\in E^\ast$ and $x_1,\ldots, x_p\in E$ 
	\begin{equation*}
		(l_1\tensor\ldots\tensor l_{k+p})\cdot (x_1\tensor\ldots\tensor x_p) = l_1\tensor\ldots \tensor l_k \gen{l_{k+1},x_1}\ldots \gen{l_{k+p},x_p}.
	\end{equation*}
	In particular, if $l\in E^\ast$, we simply denote $l^n= l\tensor\ldots\tensor l$.\\
	\item If $f$ is a differentiable map between two open sets of $E$ and $F$, $f'(x)$ is considered as a linear map belonging to $F\tensor E^{\ast}$,  $f'(x): \zeta\mapsto f'(x)\cdot\zeta$; the corresponding norm will be the standard operator norm \[\abs{f'(x)} = \sup_{\zeta\in E, \abs{\zeta}_E=1}\abs{f'(x)\cdot\zeta}_F.\]
\end{itemize}

\subsection*{Notations} In the course of the paper we will define subsets of the three spaces of germs of transformations $\G^\eta$, $\G^{\omega}$, by adding superscripts $\ast^\sigma$ and subscripts $\ast_s$. A superscript $\sigma>0$ indicates the set of $\sigma$-close to the identity diffeomorphisms, a subscript  $s>0$  indicates the $s$-width of analyticity of the diffeomorphism in the considered set.

\subsection{Space of affine non-symplectic conjugacies} 
\label{section conjugacies}
Let $$\chi_s :=\set{v\in\A(\T^n_s,\C^n): v(0)=0}$$ be the space of vector fields on the torus vanishing at the origin, endowed with the norm $\abs{\,\cdot\,}_s$. \\
Let $\mathcal{D}_s$ be the space of maps $\varphi\in\mathcal{A}(\T^n_s,\T^n_\C)$ which are real holomorphic isomorphisms from the interior of $\T^n_s$ to its image and let $\mathcal{D}^\sigma_s$ be the neighborhood of the identity, identified with maps $\varphi = \id + v:\T^n_s\to \T^n_\C$, where $v= \varphi - \id \in\chi_s$ is such that $\abs{v}_s < \sigma.$ If $\sigma$ is small enough, according to the inverse function theorem \ref{theorem well def}, such a map is a biholomorphism on $\T^n_{s'}$ for some $s'>0$.\smallskip\\
\comment{
	Let $\mathcal{D}^\sigma_s$ be the space of holomorphic invertible maps  $\varphi: \T^n_s\to\T^n_\C$, such that $\varphi(0)=0$ and \[\abs{\varphi - \id}_s<\sigma,\] where $\varphi - \id$ is viewed as a map from $\T^n_s$ to $\C^n$.\\ }
\noindent Let $\G_s$ be the affine space passing through the identity and directed by $\set{\pa{\varphi - \id, R_0 + (R_1 - {I})\cdot r}}$, where $\varphi - \id\in\chi_s$, while 
$R_0 \in \mathcal{A}(\T^n_s,\C^n)$ and  $R_1\in\mathcal{A}(\T^n_s,\GL_n(\C))$. %are such that $\Pi_{\ker A} R_0(0)=0$ and $\Pi_{\ker [A,\cdot]}(R_1(0) - {I}) = 0$. \smallskip
\comment{\footnote{This normalization, which will guaranty the uniqueness of the transformation, differs from the one chosen by Moser, who solves $u = \pull{g}(v - \lambda)$ instead. See \cite[\S 4-5]{Moser:1967}} }\\
\noindent Let then $\G^\sigma_s$ be the neighborhood of the identity in $\G_s$, consisting of the maps $$g (\teta, r) = \pa{\varphi(\teta), R_0(\teta) + R_1(\teta)\cdot r}$$ such that \[\abs{\varphi - \id}_s < \sigma\] and \[\abs{R_0(\teta) + R_1(\teta)\cdot r - r}_s < \sigma.\]

\comment{
	We define $\mathcal{G}^\sigma_s$ as the subset of $\mathcal{A}(\text{T}^n_s,\text{T}^n_\C)$ consisting of maps of the form \[g(\teta,r)=\pa{\varphi(\teta), R_0(\teta) + R_1(\teta)\cdot r},\] where $\varphi\in\mathcal{D}^\sigma_s$, while 
	$R_0 \in \mathcal{A}(\T^n_s,\C^m)$ and  $R_1\in\mathcal{A}(\T^n_s,\Mat_m(\C))$ are such that\footnote{This normalization differs from the one chosen by Moser, who solves $u = \pull{g}(v - \lambda)$ instead. See \cite[\S 4-5]{Moser:1967}} $\Pi_{\ker A} R_0(0)=0$ and $\Pi_{\ker [A,\cdot]}(R_1(0) - \operatorname{Id}) = 0$ and satisfy
	\[\abs{R_0(\teta) + R_1(\teta)\cdot r - r}_s < \sigma.\]}

%Indicating with $g_j(\teta,r)$ the $j$-th component of $g(\teta,r)$, we put \[\abs{g}_s = \max_{1\leq j\leq 2n} (\abs{g_j(\teta,r)}_s).\]
\begin{figure}[h!]
	\begin{picture}(0,0)%
		\includegraphics{Fig3-eps-converted-to.pdf}%
	\end{picture}%
	\setlength{\unitlength}{1657sp}%
	%
	\begingroup\makeatletter\ifx\SetFigFont\undefined%
	\gdef\SetFigFont#1#2#3#4#5{%
		\reset@font\fontsize{#1}{#2pt}%
		\fontfamily{#3}\fontseries{#4}\fontshape{#5}%
		\selectfont}%
	\fi\endgroup%
	\begin{picture}(7905,1881)(1561,-5473)
		\put(5581,-4426){\makebox(0,0)[lb]{\smash{{\SetFigFont{6}{7.2}{\familydefault}{\mddefault}{\updefault}{\color[rgb]{0,0,0}$g$}%
		}}}}
		\put(1846,-3751){\makebox(0,0)[lb]{\smash{{\SetFigFont{6}{7.2}{\familydefault}{\mddefault}{\updefault}{\color[rgb]{0,0,0}$\text{T}^n_{s+\sigma}$}%
		}}}}
		\put(1576,-4246){\makebox(0,0)[lb]{\smash{{\SetFigFont{6}{7.2}{\familydefault}{\mddefault}{\updefault}{\color[rgb]{0,.56,.56}$\text{T}^n_s$}%
		}}}}
		\put(1711,-4606){\makebox(0,0)[lb]{\smash{{\SetFigFont{6}{7.2}{\familydefault}{\mddefault}{\updefault}{\color[rgb]{0,0,0}$\text{T}^n_{0}$}%
		}}}}
		\put(9451,-4201){\makebox(0,0)[lb]{\smash{{\SetFigFont{6}{7.2}{\familydefault}{\mddefault}{\updefault}{\color[rgb]{0,.56,.56}$g(\text{T}^n_s)$}%
		}}}}
		\put(9406,-4651){\makebox(0,0)[lb]{\smash{{\SetFigFont{5}{6.0}{\familydefault}{\mddefault}{\updefault}{\color[rgb]{0,0,0}$g(\text{T}^n_0)$}%
		}}}}
	\end{picture}%
	
	\caption{Deformed complex domain}
\end{figure}
\noindent The "Lie Algebra" $\normal{T}_{\id}\G^\sigma_s$ of $\mathcal{G}^\sigma_s$, consists of maps \[\dot g (\teta,r)= \pa{\dot\varphi(\teta), \dot R_0(\teta) + \dot R_1(\teta)\cdot r}.\]
Here $\dot{g}$ lies in $\mathcal{A}(\text{T}^n_s,\C^{2n})$; more specifically $\dot{\varphi}\in\chi_s$, $\dot{R}_0\in\mathcal{A}(\T^n_s,\C^n)$ and $\dot{R}_1\in\mathcal{A}(\T^n_s,\Mat_n(\C))$. We endow this space with the norm \[\abs{\dot{g}}_s = \max_{1\leq j\leq 2n} (\abs{\dot{g}_j}_s).\] 

\subsection{Affine Symplectic cojugacies}\label{hamiltonian conjugacy}

We consider now the contragredient action of $\mathcal{D}^\sigma_s$ on $\normal{T}^n_s$, with values in $\normal{T}^n_\C$: \[\varphi(\teta,r) = (\varphi(\teta), \,^t\varphi'^{-1}(\teta)\cdot r).\] 
Let  $\mathcal{Z}^\sigma_s$ be the space of closed $1$-forms on $\T^n_s$ $\rho(\teta)= dS(\teta) + \xi$, where $S: \T^n_s\to \C $ and $\xi\in\R^n$ are uniquely determined assuming $S(0)=0$, such that $$\abs{\rho}_s := \max (\abs{\xi}, \abs{dS}_s)<\sigma.$$ Let $\G^{\omega,\sigma}_{s}= \mathcal{D}^\sigma_s\times \mathcal{Z}^\sigma_s$, whose elements $h = (\varphi, \rho)$ define symplectic transformations 
\begin{equation}
	h(\teta,r)= (\varphi(\teta),\,^t\varphi'^{-1}(\teta)\cdot (r + dS(\theta) + \xi)),
	\label{symplectic g}
\end{equation}

identified, locally in the neighborhood of the identity, to an open set of the affine space passing through the identity and directed by $\set{(\varphi - \id), \rho}$. 
The tangent space at the identity $T_{\id}\G^{\omega}_s=\chi_s\times\mathcal{Z}_s$,
 is endowed with the norm \[|{\dot h}|_s = \max(\abs{\dot\varphi}_s,\abs{\dot\rho}_s).\] 

\medskip

In this way we set $\cG_s := \cG^{\omega,\sigma}_s\times \cG^{\eta,\sigma}_s$.

\subsection{Spaces of vector fields} \label{Spaces of vector fields}

Concerning the Hamiltonian side, we recall that we are interested in those $K\in\mathcal{K}^\alpha$ of the form 
\begin{equation}
	K(\teta,r)= c + \alpha\cdot r + \frac{1}{2}Q(\teta)\cdot r^2 + O(r^3),
\end{equation} 
where $Q$ is a non degenerate quadratic form on $\T^n_s$ : $\det\frac{1}{(2\pi)^n}\int_{\T^n} Q(\teta)\,d\teta \neq 0$.\\
There exist $s_0$ and $\eps_0$ such that $\forall s>s_0$, $K^0\in\mathcal{H}_s$ and for all $H\in\mathcal{H}_{s}$ such that $\abs{H - K^0}_{s_0}<\eps_0$ one has \[\abs{\det\int_{\T^n}\frac{\partial^2 H}{\partial r^2}(\teta,0)\,\frac{d\teta}{(2\pi)^n}}\geq \frac{1}{2}\abs{\det\int_{\T^n}\frac{\partial^2 K^0}{\partial r^2}(\teta,0)\,\frac{d\teta}{(2\pi)^n}}\neq 0.\]
\noindent From now on, we assume that $s\geq s_0$ and define $$\mathcal{K}^\alpha_s = \set{K\in \mathcal{K}^\alpha_s: \abs{K - K^0}_{s_0}\leq \eps_0}.$$
We hence consider the corresponding space of vector fields with twist and call it 
\begin{equation}
	\Uh_s(\alpha) = \set{u^{K} (\teta,r)= (\alpha + \frac{1}{2}Q(\teta)\cdot r + O(r^2), O(r^2))}, 
\end{equation}
the affine subspace of $\Vh_s=\set{v^{\h}\in\Vh_s : \abs{H - K^0}_{s^0}\leq \eps_0}$.\\



\begin{itemize}[leftmargin=*]
	\item Let $\mathcal{W}_s= \Vh_s\oplus\cV^\eta_s\subset \cA({\normal{T}^n_s, \C^{2n}\oplus\C^{2n}})$, endowed with the norm \[\|{w}\|_s:= \max_{1\leq j\leq 2n}(|\vh_j|_s,|v^\eta_j|_s),\] 
		
		and $\mathcal{W}=\bigcup_s\,\mathcal{W}_s$.
	\item For $\alpha\in\R^n$ and $M\in\Mat_n(\R)$, let $\mathcal{U}_s(\alpha,M,\eta)$ be the affine subspace of $\cW_s$ consisting of vector fields defined in \eqref{good U} .
\end{itemize}

\subsection{The normal form operators $\chi$ and $\phi_{h,K}$}
 Let $\delta$ in Theorem \ref{Kolmogorov} be taken small enough, say $\delta < \min \{s, \sigma^2/2n\}$. Then, according to Theorem \ref{theorem well def} and Corollary \ref{cor well def}, the family of operators 
\begin{equation}
	\chi: \G_{s+\sigma}^{\omega,\delta}\times\G^{\eta,{\sigma/n}}_{s+\sigma}\times\U_{s+\sigma}(\alpha,M,\eta)\times\Lambda^\eta\to \cW_s,\,\quad (h,g^\eta,u,\lambda)\mapsto \push{h}\push{g}^\eta u + \lambda^\eta,
	\label{normal operator}
\end{equation} 

and, more specifically, fixing any $h,K$

\begin{equation}
\label{phi intermedio} 
\phi_{h,K}:   \cG_{s+\sigma}^{\eta,\sigma/n}\times\cU_{s+\sigma}^\eta(M)  \times \Lambda^\eta \to h^{*}\cW_s\, \quad  ( g^\eta, u^\eta, \lambda^\eta)\mapsto  g^\eta_*(X_K + u^\eta) + h^*\lambda^\eta
\end{equation}

are now defined. 
\\

Theroem \ref{bridge thm} and, consequently \ref{Main NFT}, follows from the invertibility of the above map $\phi_{h,K}$.

\smallskip

The drawback of focusing on $\phi_{h,K}$ is that we will need the germs of $ g^\eta_*(X_K + u^\eta) + h^*\lambda^\eta $ and $h^*(\vh + v^\eta)$ to match on the unknown torus $g(\normal{T}^n_0)$ and need to pay attention to composition operators in order not to shrink artificially the domains of analyticity, because of the rigidity of analytic maps. In this purpose, in general given a diffeomorphism $g\in\mathcal{G}^{\sigma,\eta}_{s}$ and a real analytic vector field $v$ on $g(\normal{T}^n_s)$, we define the following deformed norm \[\abs{v}_{g,s} := \abs{{g}^* v}_s\]  depending on $g^\eta$.




\subsection{Cohomological equations}
\red{Here are the general statements for matrices with possibly complex eigenvalues, where that Melnikov conditions are required. We can state it in our case...}

Here we recall three classical Lemmata, used for straightening the tangent and normal dynamics on the torus. \\ 
A vector $\alpha \in\R^n$ is identified with a constant vector field on the torus $\T^n$, thus with the derivation operator
\begin{equation*}
L_{\alpha} : \mathcal{A}(\T^n_{s+\sigma}) \to \mathcal {A}(\T^n_s),\quad f \mapsto L_\alpha f = f'\cdot \alpha := \sum_{j=1}^n \alpha_j \frac{\partial f}{\partial\teta_j}.
%\label{cohomological1}
\end{equation*}
Let now $\alpha\in\R^n$ and $M\in\Mat_m(\R)$, a diagonalizable matrix of simple, non zero eigenvalues $\mu = (\mu_1,\ldots,\mu_m)\in\C^m$ satisfy the following Diophantine conditions
\begin{align}
\label{Dio 1}
&\abs{k\cdot\alpha}\geq \frac{\gamma}{\abs{k}^\tau},\qquad \forall k\in\Z^n\setminus\set{0}\\
\label{Dio 2}
&\abs{\imath k\cdot\alpha + \mu_j}\geq \frac{\gamma}{\pa{1 + \abs{k}}^\tau},\qquad \forall k\in\Z^n, j=1,\dots,m,\\
\label{Dio 3}
&\abs{\imath k\cdot\alpha + l\cdot \mu}\geq \frac{\gamma}{\pa{1 + \abs{k}}^\tau},\qquad \forall (k,l)\in\Z^n\times\Z^m\setminus\set{0},\quad\abs{l}=2.
\end{align}

\begin{lemma}[Straightening the dynamics on the torus]
Let $\alpha\in\R^n$ satisfy condition \eqref{Dio 1}. For every $g\in\mathcal{A}(\T^n_{s+\sigma})$ of zero average, there exists a unique preimage $f\in\mathcal{A}(\T^n_s)$ by $L_\alpha$ of zero average  satisfying the following estimate \[\abs{f}_{s}=\abs{L_{\alpha}^{-1}g}_s\leq \frac{C_1}{\gamma}\frac{1}{ \sigma^{n+\tau}}\abs{g}_{s+\sigma},\] $C_1$ being a constant depending only on the dimension $n$ and the exponent $\tau$. 
\label{tangent lemma}
\end{lemma}


\noindent Let us define
\begin{equation*}
L_\alpha + M : \mathcal{A}(\T^n_{s+\sigma},\C^m) \to \mathcal {A}(\T^n_s,\C^m),\quad f\mapsto L_\alpha f + M\cdot f = f'\cdot \alpha + M\cdot f.
\label{cohomological2}
\end{equation*}
\begin{lemma}[Relocating the torus]
Let $\alpha\in\R^n$ and $M\in\GL_m(\R)$ satisfy the Diophantine condition \eqref{Dio 2}. For every $g\in\mathcal{A}(\T^n_{s+\sigma},\C^m)$, there exists a unique preimage $f\in\mathcal{A}(\T^n_{s},\C^m)$  by $L_\alpha + M$. Moreover the following estimate holds
\[\abs{f}_{s}=\abs{\pa{L_\alpha + M}^{-1} g }_s\leq \frac{C_2}{\gamma}\frac{1}{ \sigma^{n+\tau}}\abs{g}_{s+\sigma}, \] $C_2$ being a constant depending only on the dimension $n$ and the exponent $\tau$.
\label{normal lemma}
\end{lemma} 

\noindent Eventually, let us consider the space of analytic functions $F$ on $\T^n_{s+\sigma}$ with values in $\Mat_m(\C)$ and define the operator 
\begin{equation*}
\amat{llcl}{
L_\alpha + [M, \cdot]:\, & \A(\T^n_{s+\sigma},\Mat_m(\C)) &\to &\A(\T^n_s,\Mat_m(\C))\\
& F &\mapsto &L_\alpha F + [M,F]}.
\label{cohomological3}
\end{equation*}
With the notation $L_\alpha F$ (or $F'\cdot\alpha$) we mean  that we are applying the Lie derivative operator to each component $F^i_j$ of the matrix $F$; $[M,F]$ is the usual commutator.

\begin{lemma}[Straightening the first order dynamics] Let $\alpha\in\R^n$ and $M\in\GL_m(\R)$, satisfy the Diophantine conditions \eqref{Dio 1} and \eqref{Dio 3}. For every $G\in\A(\T^n_{s+\sigma},\Mat_m(\C))$ whose diagonal elements have zero average $\int_{\T^n} G^i_i \, \frac{d\teta}{(2\pi)^n}= 0$, there exists a unique $F\in\A(\T^n_{s},\Mat_m(\C))$ with  $\int_{\T^n} F^i_i \, \frac{d\teta}{(2\pi)^n}= 0$, such that the matrix equation \[L_\alpha F + [M,F] = G\] is satisfied; moreover the following estimate holds
\[\abs{F}_s \leq \frac{C_3}{\gamma}\frac{1}{\sigma^{n+\tau}}\abs{G}_{s+\sigma},\]
$C_3$ being a constant depending only on the dimension $n$ and the exponent $\tau$. 
\label{matrix lemma}
\end{lemma}

For the proofs we refer to \cite[Section 2.6]{Massetti:ETDS}. 

\section{{Proof of Main Normal Form Theroem \ref{Main NFT}}}







\subsection{Inversion of the operator $\phi_{h,K}$: estimates of ${\phi'}_{h,K}^{-1}$ and $\phi_{h,K}''$}

In the following we present the main result of this first part, from which Theorem \ref{bridge thm} follow, hence the main normal form Theorem \ref{Main NFT}. 

\begin{thm}\label{inverse bridge}
	 For any $(h,K)\in \G^{\omega,\delta}_{s}\times\Uh_s(\alpha)$ given by Kolmogorov Theorem \ref{Kolmogorov}, the operator $\phi_{h,K}$ is a local diffeomorphism in the sense that for any $\delta'< s <s+\sigma<1$ there exist $\eps_2>0$ and a unique $C^\infty$-map $\psi_{h,K}$ $$\psi_{h,K} : h^*\cW^{\eps_2}_{s+\sigma}\to \G^{\eta,\delta'}_{s}\times\U^\eta_s(M)\times\Lambda^\eta$$ such that $\phi_{h,K}\circ\psi_{h,K} = \id.$ Moreover $\psi_{h,K}$ is Whitney-smooth with respect to $(\alpha,M)$. 
\end{thm}

\bigskip 

\red{Note that $h, K^0$ are fixed and close to the identity and $K^0$ respectively, and $\phi_{\id,K^0}$ depends smoothly on them, so they do not play any role in the scheme so it suffices to prove the invertibility of $\phi_0 = \phi_{\id,K^0}$. (Bassam: dire meglio, per ora lo lascio con $h^*$ e alleggerisco notazioni}

\bigskip

  


\medskip




 
\begin{rmk}
Note that $g^\eta$ being affine in $\eta$ and close to the identity, the purely hamiltonian terms $\vh$ $X_K$ and $h$ will remain fixed, and only $O(\eta)$-kind of terms would undergo a variation through the linearization scheme and must be redressed in the inversion procedure. 
\end{rmk}
\red{It suffices to check the invertibility of $\phi_0 := \phi'_{\id,K^0}$}


Let $\overrightarrow{\U}^\eta$ be the vector space directing $\cU^\eta(M)$ and let 

\[\phi_{h,K}'(g^\eta,u^\eta,\lambda^\eta): T_{g^\eta}\mathcal{G}^{\eta,\sigma/n}_{s+\sigma}\times\overrightarrow{\mathcal{U}}^\eta_{s+\sigma}\times\Lambda^\eta\to (h^{*}\mathcal{V}^\eta)_{g^\eta,s}\] if $g^\eta$ is close to the identity, by solving 
\begin{equation}\phi_{h,K}'(g^\eta,u^\eta,\lambda^\eta)\cdot (\delta g^\eta,\delta u^\eta, \delta \lambda^\eta) = h^{*}\delta v^\eta.
\label{linearized equation}
\end{equation}


\begin{prop}[Inversion of $\phi'_{h,K}$]
\label{bound dphi-1}
If $g^\eta$ is close enough to the identity, for any $h^{*}\delta v^\eta$ in $h^{*}\mathcal{V}^\eta_{g^\eta,s+\sigma}$ there exists a unique triplet $(\delta g^\eta,\delta u^\eta, \delta \lambda^\eta)$ \comment{$\in T_g\mathcal{G}_{s}\times\overrightarrow{\mathcal{U}_{s}}\times\Lambda$} such that equation \eqref{linearized equation} is satisfied.\\ 
Moreover, there exist $\tau', C' > 0$ such that
\begin{equation}
\max\,(\abs{\delta g^\eta}_s,\abs{\delta u^\eta}_s,\abs{\delta\lambda^\eta})\leq \frac{C'}{\sigma^{\tau'}}\abs{h^{*}\delta v^\eta}_{g,s+\sigma},
\label{linear bounds}
\end{equation}
 where $C'$ is a constant that depends only on $n$, $\tau,\gamma$ and $\abs{\pa{g^\eta - \id, X_K + u^\eta - (\alpha, \eta M r)}}_{s+\sigma}$, and both sides of the inequality are of order $O(\eta)$.
\end{prop}

\begin{proof}
Let us differentiate $\phi_{h,K}$ at $(g^\eta,u^\eta,\lambda^\eta)$, where $X_{K}, h$, and $\vh$ are given. 
\medskip



%\gr{Notations}: $$u = \uh + \hu = X_K + V,\quad \tf(\teta,r) := h\circ g(\teta,r) = (\mathtt{\Theta}(\teta), \mathtt{R}_0(\teta) + \mathtt{R_1}(\teta)\cdot r),$$ where $\tf$ is affine and each component consists of Hamiltonian and non Hamiltonian $O(\eta)$ terms.\\ 
%
%\begin{equation}
%\phi: (g,h,u,\Lambda)\mapsto \push{\tf}u + \Lambda.
%\end{equation}

Since,  for any diffeomorphism $G$, the Fr\'echet derivative in the direction $\delta G$ of the push-forward operator $\push{G}: X\mapsto \push{G}X$ at $G$ is $$D(\push{G}X)\cdot\delta G = [\push{G}X, \delta G\circ G^{-1}],$$  the differentiation of $\phi_{h,K},$ w.r.t.  $ g^\eta,  u^\eta, \lambda^\eta$ in the direction of $\delta g^\eta, \delta u^\eta, \delta \lambda^\eta$ yields

%in the case of $\tf$ we have 
%\begin{equation} \label{push diff}
%[g_*h_*u, \delta \tf\circ \tf^{-1}] = [h_*g_*u, \delta h\circ h^{-1}  + \pa{h'\cdot(\delta g\circ g^{-1})}\circ h^{-1}],
%\end{equation}
%
%where we used standard rules on the Fréchet differentials over compositions. 
%Note that, coherently, expression \eqref{push diff} defines a vector field in a neigborhood of $h\circ g(\T^n\times\set{0})$.\\
%
%
%Now, we want to solve, for any fixed $x = (g,h,U,\Lambda)$, given $\delta y$ close to $\delta U =  \delta\uh + \delta\hu$, the equation $\phi(x)\delta x = \delta y$. Namely

%\begin{equation}\label{eq:linear}
% [h_*g_*u, \delta h\circ g^{-1}]  + [h_*g_*u, [h'\cdot(\delta g\circ g^{-1})]\circ h^{-1}] + h_*g_*\delta u + \delta\Lambda = \delta y.
%\end{equation} 

%In order to do so we have to pull equation \eqref{eq:linear} back by $h\circ g$ the equation  and perform a Hamiltonian step - which is classic - and a non-Hamiltonian one - which is similar to Moser's general step (see Massetti ETDS) - in order to match both terms in the equation, the Hamiltonian and the dissipative one respectively. 
%
%\begin{rmk}
%Note that because of the affine structure of the diffeomorphism $h\circ g = \id + h + \eta \hat g$ (recall also that both $h$ and $g$ are isotopic to the identity), where $\hat g$ contains mixed terms involving both $h$ and the non-hamiltonian ones, the linearization of $\phi$ will split in purely Hamiltonian terms which shall be determined first, and mixed terms (involving also the already determined Hamiltonian ones) of order $O(\eta)$.
%\end{rmk}

%By naturality of the Lie brackets and the linearity of the pull-back operator, the action of $h^\ast$  yields
%\begin{equation}
%\sq{g_*u, h'^{-1}\cdot\delta h } + \sq{ g_*u, \delta g\circ g^{-1}} + g_*\delta u + h^*\delta\Lambda = h^*\delta y,
%\end{equation}
%and the one of $g^*$

\begin{equation}
 \sq{ g^\eta_*(X_K + u^\eta), \delta g^\eta\circ g^{\eta -1}} + g^\eta_*\delta u^\eta + h^*\delta\lambda^\eta = h^*\delta v^\eta,
\end{equation}

which, by naturality of the Lie brackets, we can pull it  back by $g^\eta$ to 


\begin{equation*}
   [X_K + u^\eta, {g^\eta}'^{-1} \delta g^\eta)] +  \delta u^\eta + {g^\eta}*h^*\delta\lambda^\eta = {g^\eta}^*h^* \delta v^\eta
\end{equation*}

and obtain  an equation between vector fields at $\T^n\times\set{0}$ - as opposed to the unknown $h\circ g(\To)$.

{For ease of notations we set \[\dot\lambda^\eta:=\pull{g{^\eta}}\pull{h} \delta \lambda^\eta,\qquad \dot v^\eta:=\pull{g{^\eta}}\pull{h}\delta v^\eta, \qquad \dot{g}^\eta:= g{^\eta}'^{-1}\cdot\delta g^\eta,\]}
 and rewrite 
 
 \begin{equation}\label{pullback eq}
   [X_K + u^\eta,  \dot g^\eta] +  \delta u^\eta + \dot\lambda^\eta= \dot v^\eta
\end{equation}
 
where, in particular,
$$\dot g^\eta = (\dot{\mathtt \Theta}^\eta, \dot\xi^\eta + \dot{\mathtt R}^\eta_{0} +  \dot{\mathtt R}^\eta_{1}\cdot r  ) :=  \pa{\heta \dot\varphi (\teta),{\eta}(\dot\xi + \dot R_0(\teta) + r\cdot\dot R_1(\teta)}.$$ 

and we denoted $\dot\xi^\eta\in\R^{n}$ the constant part of $\dot{\mathtt R}^\eta_{0} $ - to be determined as well. 
We are interested in normalizing the dynamics tangentially at order zero with respect to
$r$, and up to the first order in the normal direction. To this purpose,  note that $\delta u^\eta = \eta(O(r), O(r^{2}))$ will be determined later by identification of the remainders.

\medskip


%where  $\dot h = (\dot\varphi(\teta), - \dot\varphi(\teta)\cdot r + \dot\rho(\teta) + \dot\xi )$ and $\dot g = (\dot\hp(\teta), \dot R_0(\teta) + \dot R_1(\teta)\cdot r)$



%Now, since $\delta y = \delta v^h + \heta \delta v^\eta$  hand both $h$ and $g$ are close to the identity, by linearity of the pull-back and by Taylor expanding at the identity, the r.h.s. reads 
%
%\begin{equation}
%h^*\delta\vh +  h^*[\delta\vh, g - \id] + h^*O(|g - \id|^2, \delta v^{\h}) + \heta h^*g^* \delta\hv =: \dot{v}^H+ \heta\dot{v}^\eta,
%\end{equation}
%and consist in the sum of a purely Hamiltonian vector field and a non-Hamiltonian one which is of order $O(\heta)$ and, of course, depends on the action of $h$ also. 
%The idea is to match  first  the Hamiltonian terms  of  equation \eqref{eq:linear}, then the non-Hamiltonian ones.


%Since $U = X_K + \heta V$, by linearity of the Lie brackets, we get 
%
%\begin{equation}\label{linear eq}
%	[X_K,\dot h] + [V,\dot h] + [X_K,\eta\dot{g}] + [V,\eta\dot g] + \delta U + \dot\Lambda = \dot v^{\h} + \dot v^\eta(g,h),
%\end{equation}
%where $\dot\Lambda = h^*\delta\Lambda + O(\heta,g-\id,\delta\Lambda)= (0, \dot b + \heta\dot B\cdot r) + O(\heta,g-\id,\delta\Lambda)$, with $\dot B = \delta B $ and $\dot b = (\delta b\cdot \varphi'^{-1}) + \heta\delta B(\rho\circ\varphi^{-1} + \xi)\cdot \varphi'^{-1}$. 


By expanding the vector fields along $\normal{T}^n_{0}$
\begin{equation}
\eqsys{
(X_K + u^\eta) (\teta,r)= \pa{\alpha + \mathtt{u}_1(\teta)\cdot r + O(r^2), \eta M\cdot r + \mathtt{U}_2(\teta)\cdot r^2 + O(r^3)}\\
%\dot {\tf}(\teta,r) = \pa{\dot{\mathtt{\Theta}}(\teta) ,  \dot{\mathtt{R}}_0(\teta)+ \dot{\mathtt{R}}_1(\teta)\cdot r} \\
\dot\lambda^\eta(\teta,r) = \pa{0 , \dot \Lambda^\eta_0(\teta) + \dot \Lambda^\eta_1(\teta)\cdot r }\\
\dot v^\eta(\teta,r) = \pa{\dot{\mathtt{v}}^\eta_0(\teta) + \eta O(r), \dot{\mathtt{V}}_0^\eta(\teta) + \dot{\mathtt{V}}^\eta_1(\teta)\cdot r + \eta O(r^2)},}
\label{sviluppi}
 \end{equation}
 
 where \\
 - $\mathtt{u}_1(\teta) = Q(\teta) + \eta u_1(\teta),\quad
  \mathtt{U}_2(\teta) = -\frac12 Q''(\teta) + \eta U_2(\teta)$ \\
  - $\dot v^\eta$ is of order $O(\eta)$\\
  - $\dot\Lambda^\eta$ is affine in $\delta\lambda^\eta$,  of the form $[(\id + O(|h - \id|) + \eta O(|g^\eta - \id|)) \delta\lambda^\eta$, hence of order $O(\eta)$ as well.

 
We shall match the germs of equation \eqref{pullback eq} up to order  $0$ and $1$ in $r$ in the $\teta$ and $r$ directions respectively. That is 

\begin{align}
\label{Teta} 
\dot{\mathtt{\Theta}^\eta}' \cdot \alpha - \mathtt{u}_1\cdot (\dot\xi^\eta + \dot{\mathtt{R}}^\eta_0) &= \dot {\mathtt{v}}^\eta_0 ,\\
\label{R_0}
\dot{\mathtt{R}^\eta}_0'\cdot \alpha - \eta M\cdot (\dot\xi^\eta + \dot{\mathtt{R}}^\eta_0) &= \dot{\mathtt{V}}^\eta_0 - \dot{\Lambda}^\eta_0,\\
\label{R_1}
 \dot{\mathtt{R}}_1^{\eta'}\cdot\alpha + [\eta M,\dot{\mathtt{R}}^\eta_1] + \dot{\mathtt{R}}_0^{\eta'}\cdot \mathtt{u}_1 - 2 \mathtt{U}_2\cdot \dot{\mathtt{R}}^\eta_0 &= \dot{\mathtt{V}}^\eta_1 - \dot{\Lambda}^\eta_1.
\end{align}

which, by definition of the terms given above and highlighting the dependence on $\eta$ it yields 
	\begin{align}
			\tag{$\partial_\teta$ order 0}
		\label{zero teta}
		 \heta \dot\varphi'\cdot\alpha - \pa{Q(\teta) + \heta u_1} \cdot(\eta\dot\xi  + \heta\dot{R}_0) &=  \heta\dot v_0 ,\\
		\tag{$\partial_r$ order 0}
		\label{zero r}
		{\eta}\dot R_0'\cdot\alpha+ \eta M\cdot (\heta \dot R_0 + \eta\dot\xi)  &=  {\eta}\dot{V}_0 - \eta\dot{\Lambda}_0,\\
		\tag{$\partial_r$ order 1}\label{order 1 r}
		  \heta\dot R_1'\cdot\alpha + [\heta \dot R_1,\heta M] +  \heta\dot R_0'\cdot ( \heta u_1 + Q) \\
	\notag
	 - 2\mathtt{U}_2\cdot ( \eta\dot\xi + \eta\dot R_0) &=  {\eta}\dot{V}_1 - \eta\dot{\Lambda}_1,
	\end{align}
	
	
%		\begin{align}
%			\tag{$\partial_\teta$ order 0}
%		\label{zero teta}
%		\dot{\varphi}'\cdot\alpha + \heta \dot\varphi_{\heta}'\cdot\alpha - \pa{Q(\teta) + \heta u_1} \cdot(\dot\rho + \dot\xi  + \heta\dot{R}_0) &= \dot v_0^{H} + \heta\dot v_0 ,\\
%		\tag{$\partial_r$ order 0}
%		\dot\rho'\cdot \alpha +\hat{\eta}\dot R_0'\cdot\alpha+ \eta M\cdot (\heta \dot R_0 + \dot\rho + \dot\xi)  &= \dot{V}_0^{H} + \hat{\eta}\dot{V}_0 - \dot{\Lambda}_0,\\
%		\tag{$\partial_r$ order 1}\label{order 0 r}
%		\underbrace{- \,^t D\dot\varphi' \cdot \alpha + \,^t D(Q\cdot(\dot\rho + \dot\xi))}_{\text{pure Hamiltonian terms}} + \heta\dot R_1'\cdot\alpha + [\heta \dot R_1,\heta M] 	+ [\,^t\dot\varphi', \heta M] 
%		\\
%		\notag
%	+ (\dot\rho + \dot\xi + \heta\dot R_0)'\cdot  \heta u_1 + 2 Q''\cdot  \eta\dot R_0  - 2\eta U_2\cdot (\dot\rho + \dot\xi + \eta\dot R_0) &= \dot{V}_1^{H} + \hat{\eta}\dot{V}_1 - \dot{\Lambda}_1,
%	\end{align}

It is a triangular system which, starting from equation \eqref{zero r}, determines the unknows. In the following we do not keep track of pure constants, which may depend also on the dimension, which is fixed, and denote inequalities as $\lesssim$. We shall repeatedly make use of Cauchy and triangular inequalities.

 Note that in order to solve the first equation we shall make use of the nondegenerate twist $Q$ and determine $\dot\xi$, while we should make use of the counter-terms to solve the Homological equation in the direction of the actions.\\
 We shall find solutions ``modulo $\dot\Lambda$``, which will be determined later, to absorb the average of the r.h.s. of the equations in the $r$-direction. Starting from the second equation, let 
 
 $$
 \bar b^\eta := \eta \int_{\T^n} \dot V_0 - \dot\Lambda_0\,\frac{d\teta}{(2\pi)^n} -\eta\dot\xi
 $$
 
 we get 

\begin{equation}
\eta\dot R_{0} = (L_{\alpha} + \eta M)^{{-1}} ({\eta}\dot{V}_0  -\eta \dot{\Lambda}_0 - \bar b^\eta )
\end{equation}

with

\[\abs{\eta \dot{R}_0}_{s} \leq \frac{C}{\gamma}\,\frac{|\eta|}{\sigma^{n+\tau}}\abs{\dot{V}_0 - \dot\Lambda_0 }_{s+\sigma}. \]

Note that the constant part being handled by $\Lambda_{0}$, the small denominators $|\im\alpha\cdot k + \eta\mu_{j}|$ can be bounded uniformly w.r.t. $\eta$ and the standard Diophantine condition on $\alpha$. Moreover, everything is smooth in $\eta$ and the equation is satisfied for \emph{any} $\eta \in [-\eta_{0},\eta_{0}]$.

\medskip

Taking the average of equation \eqref{zero teta}, we determine $\dot\xi$:

\begin{equation}
\eta\dot\xi = - \pa{\int_{\T^{n}} Q + \eta u_{1}} ^{-1}\cdot\pa{  \int_{\T^{n}} \eta \dot v_{0} - \eta^{2} u_{1}\cdot \dot R_{0} - \eta Q\cdot\dot R_{0}} \, d\teta
\end{equation}
\red{Note that $Q$ is nondegenerate, $\eta$ fixed and $u_{1}$ will be in a $\delta$-neighborhood of the initial datum $x^0 = (\id, u^{0,\eta},0)$ in the inverse function theorem (see Theorem \ref{teo inversa}), the nondegeneracy is preserved and the inverse well defined.}

With 
$$
|\eta \dot\xi| \lesssim |\eta| |\dot v_{0} - \eta u_{1}\cdot \dot R_{0} -  Q\cdot\dot R_{0}|
\lesssim \frac{1}{\gamma\sigma^{2(\tau + n)}} |h^*\delta v^\eta|_{g^\eta,s+\sigma}  $$
The average free part determines $\eta\delta\varphi$ with
$$
\abs{\eta\dot\varphi}_{s-\sigma} \lesssim \frac{1}{\gamma^{2}\sigma^{2\tau + 2n}} \abs{\dot v^\eta}_{s+\sigma}
$$

Thirdly, the $\Mat_m(\C)$-valued solution of \eqref{R_1} reads
\[\eta \dot{R}_1 = (L_\alpha + [\eta M,\cdot])^{-1} (\eta \dot{\tilde{V}}_1 - \eta\dot{\Lambda}_1 - \bar{B}^\eta),\]
having defined $\eta\dot{\tilde{V}}_1= \eta\dot{V}_1 -\eta \dot{R}_0'\cdot (\eta u_1 + Q) +2 \eta U_2\cdot  (\eta\dot\xi + \eta\dot{R}_0$), and 
$$\bar{B}^\eta=\int_{\T^n} \eta \dot{\tilde{V}}_1- \eta\dot{\Lambda}_1 \,\frac{d\teta}{(2\pi)^n}.$$

It now remains to handle the choice of $\delta\lambda^\eta$ such that $$\bar\lambda^\eta(\teta,r) := (0,\bar b^\eta + \bar B^\eta\cdot r)=0.$$
Note that $\bar\lambda^\eta$ belongs to $\Lambda^\eta$. Let us define a map $$ F_{h,g}: \Lambda^\eta\to \Lambda^\eta,\quad \delta\lambda^\eta\mapsto - \bar\lambda^\eta $$ in the neighborhood of $\delta\lambda^\eta=0$. %associating to $\delta\lambda$ the constant part of the known terms which cannot be solved in the cohomological equations.
It is affine and, when $g^\eta$ and $h$ are sufficiently close to the identity, invertible. \footnote{More specifically, the system in $\Lambda$ that solves $\bar\lambda^\eta =0$ is a linear system of $N$ equations in $N$ unknowns $(0,  \delta b^\eta, \delta B^\eta)$ 
%of the form $\average{a(g,\dot{v}) + A(g)\cdot\delta\lambda}=0$,
with diagonal close to $1$ if $g,h$ are close to the identity.}
\comment{ say $\eps_0$-close to the identity in the $\abs{\,\cdot\,}_{s_0}$-norm ($s_0<s<s+\sigma$),} 
\comment{$F'$ is bounded away from $0$, and $\delta\lambda\mapsto -\bar \lambda$ is invertible. Note in particular that $-\bar\lambda$ is affine in $\delta\lambda$, the system to solve being triangular of the form $\average{a(g,\dot{v}) + A(g)\cdot\delta\lambda}=0$, with diagonal close to $1$ if the smalleness condition above is assumed,}

Hence, there exists a unique $\delta\lambda^\eta$ ( Note also that $\delta\lambda^\eta = O(\eta)$) such that $F_{h,g}(\delta\lambda^\eta)=0$, satisfying $$\abs{\delta\lambda^\eta} \leq \frac{C}{\gamma\sigma^{\tau + n +1}}\abs{\dot v^\eta}_{s+\sigma}.$$

We finally have  
\[\abs{\dot{g}^\eta}_{s-2\sigma} \leq \frac{C}{\gamma^2\sigma^{2(\tau +n) +1}}\abs{h^*\delta v^\eta}_{g^\eta,s+\sigma}. \]
By definition of $\dot g^\eta$, we have $\delta g^\eta = {g^\eta}'\cdot \dot g^\eta$. In order to uniquely determine $\delta g^\eta$ we shall fix its constant term so to meet the conditions $\delta\varphi(0)=0, \delta R_0(0)=0$ and $\delta R_1(0)=0$. Since $g^\eta$ is close to the identity, these equations have a unique solution and similar kind of estimates hold for $\delta g^\eta$:

\[\abs{\delta g^\eta}_{s-2\sigma} \leq \sigma^{-1}(1 + \abs{g^\eta - \id}_{s+\sigma}) \frac{C}{\gamma^2\sigma^{2(\tau +n) +1}}\abs{h^*\delta v}_{g^\eta,s+\sigma}. \] 
Eventually, we see that $\delta u^\eta$ is actually well defined in $\overrightarrow{\U}^\eta_{s-3\sigma}$ and have \[\abs{\delta u^\eta}_{s-3\sigma}\leq \frac{C}{\gamma^2\sigma^{2(\tau + n) +3}}\abs{h^*\delta v}_{g^\eta,s+\sigma}.\]
Letting $\tau'= 2(\tau + n) + 3$ and $C'= C/\gamma^2$, up to defining $\sigma' = \sigma/4$ and $s'= s+\sigma$, the proposition is proved for all indexes $s'$ and $\sigma'$ with $s'<s' + \sigma'$.

\end{proof}
\begin{lemma}[Bounding $\phi_{h,K} ''$]
\label{bound phi''}
	The bilinear map
	\[\phi_{h,K} ''(x^\eta) : (T_g\mathcal{G}_{s+\sigma}^{\sigma/n}\times\overrightarrow{\mathcal{U}}^\eta_{s+\sigma}\times\Lambda^\eta)^{\otimes 2}\to h^*\mathcal{V}^\eta_{s}, \] where $x^\eta = (g^\eta,u^\eta,\lambda^\eta)$,
	%\[\sq{\sq{\push{g}u,\delta g\circ g^{-1}},\delta g\circ g^{-1}} + 2 \sq{\push{g} \delta u,\delta g\circ g^{-1}} - \sq{\push{g}u, (\delta g'\cdot g'^{-1}\cdot\delta g)\circ g^{-1}}. \] 
	satisfies the following estimate \[\abs{\phi''(x^\eta)\cdot \delta x^{\eta,\otimes 2}}_{g^\eta,s}\leq \frac{C'' }{\sigma^{\tau ''}} \abs{\delta x^\eta}^{2}_{s+\sigma},\]
	$C''$ being a constant depending on $\abs{x^\eta}_s$.
	\label{lemma derivata seconda}
\end{lemma}

\begin{proof}
	 In ordet to lighten notations, let us drop the apex $\eta$. Recall the expression of 
	$$\phi'_{h,K}(x)\cdot\delta x=  h_* \sq{ g_*(X_K + \eta u), \delta g\circ g^{-1}} + h_*g_*\eta\delta u + \delta\Lambda $$ 
	Differentiating again with respect to $x$ yelds
	\[h_*\left(\sq{\sq{\push{g}(X_K + u),\delta g\circ g^{-1}} + \push{g}\delta u, \delta g\circ g^{-1}}  - \sq{\push{g}(X_K + u), \delta g'\circ g^{-1}\cdot\delta g^{-1}} + \sq{\push{g}\delta u,\delta g\circ g^{-1}}\right) \]
	Since $\delta g^{-1} = - (g'^{-1}\cdot \delta g)\circ g^{-1}$, 
	\[\pull{g}\pull{h}\phi''(x)\cdot \delta x^{\otimes 2} = 2\sq{\delta u, \dot g} + \sq{\sq{u,\dot g}, \dot g} + \sq{X_K + u,\pull{g} (\delta g'\cdot \inderiv{g}\cdot \delta g)\circ g^{-1}}, \] where the last term simplifies in \[\sq{X_K + u,\inderiv{g}\cdot (\delta g'\cdot \inderiv{g}\cdot \delta g)}.\] The wanted bound follows from repeatedly applying Cauchy's inequality, triangular inequality and Lemma \ref{lemma lie brackets}.
	
\end{proof}

\subsection{Proof of the Bridge normal form Theorem }

Bridge normal form Theorem \ref{bridge thm} follows from Theorem \ref{inverse bridge}, by application of the abstract inverse function theorem \ref{teo inversa} in Appendix \label{section teorema inversa}. More specifically, Proposition \ref{bound dphi-1} and Lemma \ref{bound phi''} guarantee to apply Theorem \ref{teo inversa} in a neighborhood of $x_0 = (\id, u^{0,\eta},0)$, provided that $\|v^\eta - u^{0,\eta}\|_{s+\sigma} < \eps_2 :=\min\{\eps_1,\eps_0\}$ where $\eps_1$ is the radius of the domain of validity of Kolmogorov Theorem, given in Theorem \ref{Kolmogorov}, while $\eps_0:=\eps$ is the threshold given in Theorem \ref{teo inversa}.
\subsection{Proof of Theorems \ref{Main NFT} and \ref{thm NF parameters}}

In order to get Theorem \ref{Main NFT}, it suffices to combine  Theorems \ref{Kolmogorov}-\ref{inverse bridge}. In fact, for any given $w = (\vh, v^\eta)\in\cW_{s+\sigma}$ such that $\| w - (u_{K^0}, u^{0,\eta})\|_{s+\sigma} < \min\set{\eps_1,\eps_2}$, there exists $h,g^\eta,\lambda^\eta, u_K, u^\eta$ respectively given by $(h,X_K) = \psi_{\tK}(\vh)$ and $(g^\eta,u^\eta,\lambda^\eta) = \psi_{h,K}(h^*v^\eta)$ so that

$$h^*(\vh + v^\eta) = g^\eta_*(X_K + u^\eta) + h^*\lambda^\eta \quad \quad \Leftrightarrow\quad \quad (\vh + v^\eta) = h_*g^\eta_*(X_K + u^\eta) + \lambda^\eta $$ 

is satisfied. 

\smallskip

Concerning Theorem \ref{thm NF parameters}, it suffices to reason on a trivial perturbation of the vectorfield $u^0 + \eta\nu\partial_r = (\alpha + Q\cdot r, \eta M r + \eta\nu)$ as follows. Let us introduce a diophantine $\alpha'$ close to $\alpha$ and a matrix $M'$ close to $M$, and rewrite $u^0 = (\alpha' + (\alpha - \alpha') + Q\cdot r, \eta Mr + \eta\nu)$. By 


\appendix
\section{Inverse function theorem \& regularity of $\phi$ }
\label{section teorema inversa}

The following results are proved in \cite{Massetti:ETDS}. For convenience of the reader and not to cause confusion with the dissipation parameter $\eta$, with respect to the corresponding statements in \cite{Massetti:ETDS}, we replaced the radius $\eta$ with $\delta$ in the statements below. 

Let $E=(E_s)_{0<s<1}$ and $F=(F_s)_{0<s<1}$ be two decreasing families of Banach spaces with increasing norms $\abs{\cdot}_s$ and let $B^E_s(\sigma)=\set{x\in E : \abs{x}_s<\sigma}$ be the ball of radius $\sigma$ centered at $0$ in $E_s$. \\ On account of composition operators, we additionally endow $F$ with some deformed norms which depend on $x\in B^E_s(s)$ such that \[ \abs{y}_{0,s} = \abs{y}_s \qquad \text{and}\qquad \abs{y}_{\hat{x},s}\leq \abs{y}_{x, s + \abs{x - \hat{x}}_s}.\]
Consider then operators commuting with inclusions $\phi: B^E_{s+\sigma}(\sigma) \to F_s$, with $0<s<s+\sigma<1$, such that $\phi(0) = 0$. \\We then suppose that if $x\in B^E_{s+\sigma}(\sigma)$ then $\phi'(x):E_{s+\sigma}\to F_s$ has a right inverse $\phi'^{-1}(x): F_{s+\sigma}\to E_s$ (for the particular operators $\phi$ of this work, $\phi'$ is both left and right invertible).\\ $\phi$ is supposed to be at least twice differentiable.\\
Let $\tau:=\tau'+\tau''$ and $C:=C'C''$.
\begin{thm}\label{teo inversa}
Under the previous assumptions, assume
\begin{align}
\label{bound phi'^-1}
\abs{\inderiv{\phi}(x)\cdot\delta y}_s\, &\leq \frac{C'}{\sigma^{\tau'}}\abs{\delta y}_{x,s+\sigma}\\
\label{bound phi''}
\abs{\phi''(x)\cdot \delta x^{\tensor 2}}_{x,s}\, &\leq \frac{C''}{\sigma^{\tau''}}\abs{\delta x}^2_{s+\sigma},\quad \forall s,\sigma: 0<s<s+\sigma<1
\end{align}
$C'$ and $C''$ depending on $\abs{x}_{s+\sigma}$, $\tau',\tau''\geq 1$. \\
For any $s, \sigma, \delta$ with $\delta<s$ and $\eps\leq \delta \, \frac{\sigma^{2\tau}}{2^{8\tau}C^2}$ ($ C\geq 1,\sigma< 3 C$), $\phi$ has a right inverse $\psi: B^F_{s+\sigma}(\eps)\to B^E_{s}(\delta)$. In other words, $\phi$ is locally surjective: $$ B^{F}_{s+\sigma}(\eps)\subset \phi(B^{E}_{s}(\delta)).$$ 
\end{thm}
Define
\begin{equation}
Q : B^E_{s+2\sigma}(\sigma) \times B^E_{s+2\sigma}\to F_s,\quad (x,\hat{x})\mapsto \phi(\hat{x}) - \phi(x) - \phi'(x)(\hat{x} - x),
\label{resto Taylor}
\end{equation}  
 the reminder of the Taylor formula.\\

\begin{lemma} 
\label{lemma Q}
For every $x,\hat{x}$ such that $\abs{x-\hat{x}}_s < \sigma$, 
\begin{equation}
\abs{Q(x,\hat{x})}_{x,s}\leq \frac{C''}{2\sigma^2}\,\abs{\hat{x} - x}^2_{{s+\sigma+\abs{\hat{x} - x}}_s}.
\label{stima Q}
\end{equation}
\end{lemma}

\begin{proof}
Let $x_t = (1-t)x + t\hat{x}$, $0\leq t\leq 1$, be the segment joining $x$ to $\hat{x}$. Using Taylor's formula,  \[Q(x,\hat{x})= \int_0^1 (1-t)\phi''(x_t)(\hat{x}-x)^2\, dt,\] hence
\begin{align*}
\abs{Q(x,\hat{x})}_{x,s}&\leq \int_0^1 (1-t)\abs{\phi''(x_t)(\hat{x}-x)^2}_{x,s}\, dt \\
&\leq \int_0^1 (1-t)\abs{\phi''(x_t)(\hat{x}-x)^2}_{x_t,s + \abs{x_t - x}_{s}}\,dt\\
&\leq \int_0^1 (1-t)\,\frac{C''}{\sigma^2} \abs{(\hat{x}-x)}^2_{s +\sigma + \abs{x_t - x}_{s}}\,dt\\
&\leq \frac{C''}{2\sigma^2}\,\abs{\hat{x} - x}^2_{{s+\sigma+\abs{\hat{x} - x}}_s}.
\end{align*}

\end{proof}

\begin{proof}[Proof of theorem \ref{teo inversa}]
Let $\delta<s, \sigma $ and $\eps$ be fixed positive real numbers and let $y\in B^F_{s+\sigma}(\eps)$. We define the following map: 
\begin{equation*}
f: B^E_{s+\sigma}(\sigma)\to E_s,\quad x\mapsto x + \inderiv{\phi}(x)(y - \phi(x)).
\end{equation*}
We want to prove that, if $\eps$ is sufficiently small, there exists a sequence defined by induction by
\begin{equation*}
\eqsys{x_0 = 0\\
x_{n+1} = f(x_n) ,}
\end{equation*}
converging towards some point $x\in B^E_s(\delta)$, the preimage of $y$ by $\phi$.\\ Let us introduce two sequences
\begin{itemize}[leftmargin=*]
\item a sequence of positive real numbers $(\sigma_n)_{n\geq 0}$ such that $3\sum_n \sigma_n = \sigma$ be the total width of analyticity we will have lost at the end of the algorithm,
\item the decreasing sequence $(s_n)_{n\geq 0}$ defined inductively by $s_0 = s + \sigma$ (the starting width of analyticity), $s_{n+1}= s_n - 3\sigma_n$. Of course, $s_n\to s$ when $n\to + \infty$.
\end{itemize}  
Let $C_k = 2C \sigma_k^{-\tau} \geq 1$ for all $k \in \mathbb{N}$ and define $\zeta_n = \prod_{0\leq k \leq n} C_k^{2^{-k}} $ and $\zeta = \prod_{k\geq 0} C_k^{2^{-k}}$. We start to prove that for every $n\geq 1$, there exist $x_0,...,x_{n}$ and that $$\abs{x_{n} - x_{n-1}}_{s_n}\leq (\eps \zeta_{n-1})^{2^{n-1}}\quad \text{and} \quad \abs{x_{n}}_{s_{n}} \leq \sum_{k= 0}^{n-1} (\eps\zeta)^{2^k}.$$  From $x_{k} - x_{k-1} =  \inderiv{\phi}(x_{k-1})(y - \phi(x_{k-1}))$ we see that $y - \phi(x_k) = - Q(x_{k-1},x_k)$, which permits to write $x_{k+1} - x_k = - \inderiv{\phi}(x_k)Q(x_{k-1},x_k)$, for $k= 1,...,n$.\\ 
First, remark that $$\abs{x_1 - x_0}_{s_1} = \abs{x_1}_{s_1} \leq \frac{C'}{(3\sigma_0)^{\tau'}}\,\abs{y-\phi(x_0)}_{s_0}\leq \frac{C}{2\sigma^{\tau}_0}\,\abs{y}_{s+\sigma}\leq C_0\eps $$
and $\abs{x_1}_{s_1} \leq \zeta\eps,$ the assertion is thus true for $n=1$.\\ 
Assuming that $\abs{x_{k} - x_{k-1}}_{s_k}\leq \sigma_k$, for $k=1,...n$, from the estimate of the right inverse and the previous lemma we get 

\begin{equation*}
\abs{x_{n+1} - x_n}_{s_{n+1}}\leq \frac{C}{2\sigma_n^{\tau}}\abs{x_n - x_{n-1}}^2_{s_{n}} \leq \ldots \leq C_n C_{n-1}^2\ldots C_1^{2^{n-1}}\abs{x_1 - x_0}^{2^n}_{s_1}.
\end{equation*}
Second, observe that since $C_k\geq 1$ (see remark below), $$\abs{x_{n+1}-x_n}_{s_{n+1}}\leq C_n( \eps \zeta_{n-1})^{2^n} = (\eps\zeta_n)^{2^n} \leq (\eps\zeta)^{2^n} =\pa{\eps\prod_{k\geq 0}C_k^{2^{-k}}}^{2^n}$$ and $\abs{x_{n+1}}_{s_{n+1}} \leq \abs{x_n}_{s_n} + \abs{x_{n+1} - x_n}_{s_{n+1}} \leq  \sum_{0}^n (\eps\zeta)^{2^k}$.\\
Third, note that $$ \sum_{n\geq 0} z^{2^n} = z + z^2 + z^4 + \ldots \leq z\sum_{n\geq 0} z^n \leq 2z, $$ if $z\leq\frac{1}{2}$.

The key point is to choose $\eps$ such that $\eps \prod_{k\geq 0}C_k^{2^{-k}}\leq \frac{1}{2}$ (or any positive number $< 1$)  and $\sum_{n\geq 0}\abs{x_{n+1} - x_n}_{s_{n+1}}<\delta$, in order for the whole sequence $(x_k)$ to exist and converge in $B_s(\delta)\subset E_s$. Hence, using the definition of the $C_n$'s and the fact that $$\pa{\frac{C}{2}}^{-2^{-k}}=\pa{\frac{2}{C}}^{\pa{\frac{1}{2}}^k}\Longrightarrow\prod  \pa{\frac{2}{C}}^{\pa{\frac{1}{2}}^k} = \pa{\frac{2}{C}}^{\sum\frac{1}{2^k}}=\pa{\frac{2}{C}}^2,$$ within $\sum_k \frac{1}{2^k} = \sum_k k\frac{1}{2^{k}}=2$, we obtain as a sufficient value
\begin{equation}
\eps = \delta \frac{2}{C^2}\prod_{k\geq 0}\sigma_k^{\tau\,(\frac{1}{2})^k}.
\label{bound epsilon}
\end{equation}
Eventually, the constraint $3\sum_{n\geq 0}\sigma_n = \sigma$ gives $\sigma_k= \frac{\sigma}{6}\,\pa{\frac{1}{2}}^k$, which, plugged into \eqref{bound epsilon}, gives: 
\[\eps = \delta \, \frac{2}{C^2}\pa{\frac{\sigma}{12}}^{2\tau} > \frac{\sigma^{2\tau}\delta}{2^{8^{\tau}}C^2},\]
hence the theorem. 

A posteriori, the exponential decay we proved makes straightforward the further assumption $\abs{x_k - x_{k-1}}_{s_k}<\sigma_k$ to apply lemma \ref{lemma Q}.\\ Concerning the bounds over the constant $C$, as $\sum_k \abs{x_{k+1} - x_k}_{s_{k+1}}\leq\delta$, we see that all the $\abs{x_n}_{s_n}$ are bounded, hence the constants $C'$ and $C''$ depending on them. \\Moreover, to have all the $C_n\geq 1$, as we previously supposed, it suffices to assume $C\geq \sigma/3$.
\end{proof}

\begin{rmk}
In case the operator $\phi$ is defined only on polynomially small balls $$\phi: B^E_{s+\sigma}(c_0 \sigma^{\ell})\to F_s,\, c_0>0, \forall s,\sigma$$ the statement and the proof of theorem \ref{teo inversa} still hold, provided that $\delta$ is chosen small enough ($\delta< 2c_0(\sigma/12)^{\ell}$ suffices). \\ This is the case of the operators defined in sections \ref{dissipative Herman} and \ref{Russmann section vector fields}, where $\ell=2$. 
\label{remark polynomial}
\end{rmk}
\subsection{Local uniqueness and regularity of the normal form}
We want to show the uniqueness and some regularity properties of the right inverse $\psi$ of $\phi$, assuming the additional left invertibility of $\phi'$ (which is the case, for the particular operator $\phi'$ of interest to us).
\begin{defn}
We will say that a family of norms $(\abs{\cdot}_s)_{s>0}$ on a grading $(E_s)_{s>0}$ is \emph{log-convex} if for every $x\in E_s$ the map $ s\mapsto \log\abs{x}_s $ is convex.
\end{defn}

\begin{lemma} If $(\abs{\,\cdot\,}_s)$ is log-convex, the following inequality holds \[\abs{x}^2_{s+\sigma}\leq \abs{x}_s \abs{x}_{s+\tilde{\sigma}},\quad \forall s,\sigma,\tilde{\sigma}=\sigma(1+\frac{1}{s}).\] 
\end{lemma}

\begin{proof}
If $f: s\mapsto \log\abs{x}_s$ is convex, this inequality holds \[f\pa{\frac{s_1 + s_2}{2}}\leq \frac{f(s_1) + f(s_2)}{2}.\] Let now $x\in E_s$, then
\[\log\abs{x}_{s+\sigma}\leq\log\abs{x}_{\frac{2s + \tilde{\sigma}}{2}}\leq \frac{1}{2}\pa{\log\abs{x}_s + \log\abs{x}_{s+\tilde{\sigma}}}=\frac{1}{2}\log(\abs{x}_s\abs{x}_{s+\tilde{\sigma}}),\]hence the lemma. 
\end{proof}
Let us assume that the family of norms $(\abs{\cdot}_s)_{s>0}$ of the grading $(E_s)_{s>0}$ are log-convex. To prove the uniqueness of $\psi$ we are going to assume that $\phi'$ is also left-invertible. 

\begin{prop}[Lipschitz continuity of $\psi$] 
\label{lipschitz}
Let $\sigma<s$. If $y,\hat{y}\in B^{F}_{s+\sigma}(\eps)$ with $\eps = {3^{-4\tau}2^{-16\tau}}\frac{\sigma^{6\tau}}{4C^3}$, the following inequality holds
\[\abs{\psi(y) - \psi(\hat{y})}_s\leq L \abs{y-\hat{y}}_{x,s+\sigma},\] with $L=2C'/\sigma^{\tau'}$. In particular, $\psi$ being the unique local right inverse of $\phi$, it is also its unique left inverse.
\end{prop}

\begin{proof}
In order to get the wanted estimate we introduce an intermediate  parameter $\xi$, that will be chosen later, such tat $\delta<\xi<\sigma<s<s+\sigma$.\\ 
To lighten notations let us call $\psi(y)=: x$ and $\psi(\hat{y})=: \hat{x}$. Let also $\eps=\frac{\xi^{2\tau}\delta}{2^{8\tau}C^2}$ so that if $y,\hat{y}\in B^{F}_{s+\sigma}(\eps)$, $x,\hat{x}\in B^E_{s+\sigma -\xi}(\delta)$, by theorem \ref{teo inversa}, provided that $\delta< s + \sigma -\xi$ - to check later. In particular, we assume that any $x,\hat{x}\in B^{E}_{s + \sigma -\xi}$ satisfy $\abs{x - \hat{x}}_{s+\sigma - \xi}\leq 2\delta.$
Writing \[(x-\hat{x})=\phi'^{-1}(x)\cdot\phi(x)(x-\hat{x}),\]
and using $$\phi'(x)(x-\hat{x})= \phi(\hat{x})-\phi(\hat{x})-Q(x,\hat{x}),$$ we get \[ x-\hat{x}= \inderiv{\phi}(x)\pa{\phi(\hat{x}) - \phi(x) - Q(x,\hat{x})}.\]
Taking norms we have
 \begin{align*}
\abs{x-\hat{x}}_s  &\leq \frac{C'}{\sigma^{\tau'}}\abs{y - \hat{y}}_{x,s+\sigma} + \frac{C}{2\xi^{\tau}}\abs{x-\hat{x}}^2_{s+2\xi + \abs{x-\hat{x}}_{s + \xi}},\\
&\leq \frac{C'}{\sigma^{\tau'}}\abs{y - \hat{y}}_{x,s+\sigma} + \frac{C}{2\xi^{\tau}}\abs{x-\hat{x}}^2_{s+2\xi + 2\delta},
\end{align*} 
 by lemma \ref{lemma Q} and the fact that $\abs{x-\hat{x}}_{s+\xi}\leq\abs{x-\hat{x}}_{s+\sigma-\xi}$ (choosing $\xi$ so that $2\xi< \sigma $ too). \\ Let us define $\tilde{\sigma}=(2\xi + 2\delta)(1 + 1/s)$ and use the interpolation inequality \[\abs{x-\hat{x}}^2_{s+2\delta +2\xi}\leq\abs{x-\hat{x}}_s\abs{x-\hat{x}}_{s+\tilde{\sigma}}\] to obtain \[(1 - \frac{C}{2\xi^\tau}\abs{x-\hat{x}}_{s+\tilde{\sigma}})\abs{x-\hat{x}}_s\leq \frac{C'}{\sigma^{\tau'}}\abs{y-\hat{y}}_{x,s+\sigma}.\]
We now choose $\delta$ so small to have
\begin{itemize}[leftmargin=*]
\item $\tilde{\sigma}\leq\sigma - \xi$, which implies $\abs{x-\hat{x}}_{s+\tilde{\sigma}}\leq 2\delta$. It suffices to have $\delta\leq\frac{\sigma}{2(1+\frac{1}{s})} - \frac{3}{2}\xi$.
\item $\delta\leq\frac{\xi^{\tau}}{2C}$ in order to have $\frac{C}{2\xi^\tau}\abs{x-\hat{x}}_{s+\sigma}\leq\frac{1}{2}.$
\end{itemize}
A possible choice is $\xi = \frac{\sigma^2}{12}$ and $\delta = \pa{\frac{\sigma}{12}}^{2\tau}\frac{1}{4C},$ hence our choice of $\eps$.
\end{proof}

\begin{prop}[Smooth differentiation of $\psi$]
\label{smoothness}
 Let $\delta<s<s+\sigma$ and $\eps$ be as in proposition \ref{lipschitz}. There exists a constant $K$ such that for every $y,\hat{y}\in B_{s+\sigma}^F(\eps)$ we have \[\abs{\psi(\hat{y}) - \psi(y) - \inderiv{\phi}(\psi(y))(\hat{y} - y)}_s\leq K(\sigma) \abs{\hat{y} - y}^2_{x,s+\sigma},\] and the map $\psi': B_{s+\sigma}^F(\eps)\to L(F_{s+\sigma},E_s)$ defined locally by $\psi'(y)=\phi'^{-1}(\psi(y))$  is continuous. In particular, if $\phi: B^E_{s+\sigma}(\sigma)\to F_s$ is $C^k$, $2\leq k\leq \infty$, so is $\psi: B^F_{s+\sigma}(\eps)\to E_s$.
\end{prop}
\begin{proof}
 Let's baptize some terms
\begin{itemize}[leftmargin=*]
\item $\Delta := \psi(\hat{y}) - \psi(y) - \inderiv{\phi}(x)(\hat{y} - y) $
\item $\delta:= \hat{y} - y$, the increment
\item $\xi:= \psi(y + \delta) - \psi(y)$
\item $\Xi:= \phi(x + \xi) - \phi(x)$.
\end{itemize}
With these new notations we can see $\Delta$ as 
\begin{align*}
\Delta &= \xi - \inderiv{\phi}(x)\cdot\Xi\\
&=\inderiv{\phi}(x)(\phi'(x)\cdot\xi - \Xi)\\
& =\inderiv{\phi}(x)(\phi'(x)\xi - \phi(x+\xi) + \phi(x))\\
& = -\inderiv{\phi}(x)Q(x,x+\xi)
\end{align*}
Taking norms we have
\[\abs{\Delta}_s\leq K \abs{\hat{y} - y}^2_{x,s+\bar{\sigma}}\] by proposition \ref{lipschitz} and lemma \ref{lemma Q}, for some $\bar\sigma$ which goes to zero when $\sigma$ does, and some constant $K>0$ depending on $\sigma$ . Up to substituting $\sigma$ for $\bar{\sigma}$, we have proved the statement.\\
In addition \[\psi'(y)= \phi^{-1}(y)'= \phi'^{-1}\circ \phi^{-1}(y)= \phi'^{-1}(\psi(y)),\] the inversion of linear operators between Banach spaces being analytic, the map $y\mapsto \phi'^{-1}(\psi(y))$ has the same degree of smoothness as $\phi'$.
\end{proof}

It is sometimes convenient to extend $\psi$ to non-Diophantine characteristic frequencies $(\alpha,a)$. Whitney smoothness guarantees that such an extension exists.
Let suppose that $\phi(x)=\phi_{\nu}(x)$ depends on some parameter $\nu\in B^d$ (the unit ball of $\R^d$) and that it is $C^1$ with respect to $\nu$ and that estimates on $\phi'^{-1}_{\nu}$ and $\phi_{\nu}''$ are uniform with respect to $\nu$ over some closed subset $D$ of $\R^{d}$. 

\begin{prop}[Whitney differentiability]Let us fix $\eps, \sigma, s$ as in proposition \ref{lipschitz}. The map $\psi : D\times B^{F}_{s+\sigma}(\eps)\to B^{E}_s(\delta)$ is $C^1$-Whitney differentiable and extends to a map $\psi : \R^{d}\times B^{F}_{s+\sigma}(\eps)\to B^{E}_s(\delta) $ of class $C^1$. If $\phi$ is $C^k$, $1\leq k\leq \infty$, with respect to $\nu$, this extension is $C^k$.
\label{Whitney}
\end{prop}

\begin{proof}
Let $y\in B^{F}_{s+\sigma}(\eps)$. For $\nu,\nu + \mu\in D$, let $x_\nu = \psi_{\nu}(y)$ and $x_{\nu + \mu}=\psi_{\nu+\mu}(y)$, implying \[\phi_{\nu+\mu}(x_{\nu+\mu}) - \phi_{\nu+\mu}(x_{\nu}) = \phi_{\nu}(x_\nu) - \phi_{\nu+\mu}(x_{\nu}).\] It then follows, since $y\mapsto\psi_{\nu+\mu}(y)$ is Lipschitz, that \[\abs{x_{\nu+\mu} - x_{\nu}}_s\leq L \abs{\phi_{\nu}(x_{\nu}) - \phi_{\nu+\mu}(x_{\nu})}_{x_{\nu},s+\sigma},\] taking ${y}=\phi_{\nu +\mu}(x_{\nu}),\hat{y}=\phi_{\nu+\mu}(x_{\nu+\mu}).$ In particular since $\nu\mapsto\phi_{\nu}(x_{\nu})$ is Lipschitz, the same is for $\nu\mapsto x_{\nu}.$
Let us now expand $\phi_{\nu+\mu}(x_{\nu+\mu})=\phi(\nu+\mu, x_{\nu+\mu})$ in Taylor at $(\nu,x_{\nu}).$ We have 
\[\phi(\nu+\mu,x_{\nu+\mu}) = \phi(\nu,x_{\nu}) + D\phi(\nu,x_{\nu})\cdot (\mu, x_{\nu+\mu} - x_{\nu}) + O(\mu^2, \abs{x_{\nu+\mu} - x_{\nu}}_s^2),\] hence formally defining the derivative $\partial_{\nu} x_{\nu} := - \phi'^{-1}_{\nu}(x_{\nu})\cdot\partial_{\nu} \phi_{\nu}(x_{\nu}),$ we obtain \[x_{\nu+\mu} - x_{\nu} - \partial_{\nu} x_{\nu}\cdot\mu = \phi'^{-1}_{\nu}(x_{\nu})\cdot O(\mu^2),\] hence \[\abs{x_{\nu+\mu} - x_{\nu} - \partial_{\nu} x_{\nu}\cdot\mu}_s = O(\mu^2)\] by Lipschitz property of $\nu\mapsto x_{\nu},$ when $\mu\mapsto 0$, locally uniformly with respect to $\nu$. Hence $\nu\mapsto x_{\nu}$ is $C^1$-Whitney-smooth and by Whitney extension theorem, the claimed extension exists. Similarly if $\phi$ is $C^k$ with respect to $\nu$, $\nu\mapsto x_{\nu}$ is $C^k$-Whitney-smooth. See \cite{Abraham-Robbin:1967} for the straightforward generalization of Whitney's theorem to the case of interest to us: $\psi$ takes values in a Banach space instead of a finite dimension vector space; but note that the extension direction is of finite dimension though.
\end{proof}












\section{Inversion of a holomorphism of $ \T^n_s $ }
We present here a classical result on the inversion of holorphisms on the complex torus $\T^n_s$ that intervened to guarantee the well definition of normal form operators $\phi$.

All complex extensions of manifolds are defined at the help of the $\ell^\infty$-norm, \[\T^n_s = \set{\teta\in\T^n_\C : \abs{\teta}:= \max_{1\leq j\leq n}\abs{\Im{\teta_j}}\leq s}.\]
Let also define $\R^n_s := \R^n \times (-s,s)$ and consider the universal covering of $\T^n_s$, $p: \R^n_s\to \T^n_s$.
\begin{thm}
	\label{theorem well def} Let $v:\T^n_s\to\C^n$ be a vector field such that $\abs{v}_s < \sigma/n$. The map $\id + v:\T^n_{s-\sigma}\to\R^n_{s}$ induces a map $\varphi = \id + v : \T^n_{s-\sigma}\to\T^n_s$ which is a biholomorphism and there is a unique biholomorphism $\psi : \T^n_{s-2\sigma}\to\T^n_{s-\sigma}$ such that $\varphi\circ\psi = \id_{\T^n_{s - 2\sigma}}.$ \\
	Furthermore, the following inequalities hold: \[\abs{\psi - \id}_{s-2\sigma}\leq \abs{v}_{s-\sigma}\] and, if $\abs{v}_s<\sigma/2n$ \[\abs{\psi' - \id}_{s-2\sigma}\leq \frac{2}{\sigma}\abs{v}_{s}.\]
\end{thm}

\begin{proof}
	Let $\hat\varphi:= \id + v\circ p : \R^n_s\to\R^n_{s+\sigma}$ be the lift of $\varphi$ to $\R^n_s$.\\Let's start proving the injectivity and surjectivity of $\hat\varphi$; the same properties for $\varphi$ descend from these.
	\begin{itemize}[leftmargin=*]
		\item $\hat\varphi$ is injective as a map from $\R^n_{s-\sigma}\to\R^n_s$.\\Let $\hat\varphi(x)=\hat\varphi(x')$, from the definition of $\hat\varphi$ we have 
		\begin{align*}
			\abs{x -x'}=\abs{v\circ p(x') - v\circ p(x)} &\leq \int_0^1 \sum_{k=1}^n\abs{\partial_{x_k} \hat{v}}_{s-\sigma}\abs{x_k'-x_k}\,dt\leq \frac{n}{\sigma}\abs{v}_s\abs{x-x'}\\ & < \abs{x - x'},
		\end{align*}
		hence $x'=x$.
		\item $\hat\varphi: \R^n_{s-\sigma}\to \R^n_{s-2\sigma}\subset \hat\varphi (\R^n_{s-\sigma})$ is surjective.\\ Define, for every $y\in\R^n_{s-2\sigma}$ the map \[f: \R^n_{s-\sigma}\to\R^n_{s-\sigma},\, x\mapsto y - v\circ p(x),\] which is a contraction (see the last but one inequality of the previous step). Hence there exists a unique fixed point such that $\hat{\varphi}(x) = x + v\circ p(x) = y.$ 
	\end{itemize}
	For every $k\in 2\pi\Z^n$, the function $\R^n_s\to\R^n_s,\, x\mapsto \hat\varphi(x+k) - \hat\varphi(x)$ is continuous and $2\pi\Z^n$-valued. In particular there exists $A\in\GL_n(\Z)$ such that $\hat\varphi(x + k)= \hat\varphi(x) + Ak$.
	\begin{itemize}[leftmargin=*]
		\item $\varphi: \T^n_{s-\sigma}\to\T^n_s$ is injective.\\ 
		Let $\varphi(p(x))=\varphi(p(x')),$ with $p(x),p(x')\in\T^n_{s-\sigma}$, hence $\hat\varphi(x')=\hat\varphi(x) + k' ,$ for some $k'\in 2\pi\Z^n$. Hence $\hat\varphi(x' - A^{-1}k')=\hat\varphi(x)$, and for the injectivity of $\hat\varphi$, $p(x)=p(x')$. In particular $\varphi$ is biholomorphic:
		\begin{lemma}[\cite{Fritzsche-Grauert}] If $G\subset\C^n$ is a domain and $f: G\to\C^n$ injective and holomorphic, then $f(G)$ is a domain and $f: G\to f(G)$ is biholomorphic.
		\end{lemma}
		\item That $\varphi: \T^n_{s-\sigma}\to \T^n_{s-2\sigma}\subset \varphi(\T^n_{s-\sigma})$ is surjective follows from the one of $\hat\varphi$.
		\item Estimate for $\psi: \T^n_{s-2\sigma}\to \T^n_{s-\sigma}$ the inverse of $\varphi$.\\ Let $\hat\psi : \R^n_{s-2\sigma}\to\R^n_{s-\sigma}$ be the inverse of $\hat\varphi$, and $y\in\R^n_{s-2\sigma}.$ From the definition of $\hat\varphi$, $v\circ p(\hat\psi(y))= y - p(\hat\psi(y)) = y - \hat\psi(y)$. Hence \[\abs{\hat\psi(y) - y}_{s-2\sigma} = \abs{v\circ p (\hat \psi (y))}_{s-2\sigma}\leq \abs{v}_{s-2\sigma}\leq\abs{v}_{s-\sigma}.\]
		\item Estimate for $\psi' = \varphi'^{-1}\circ\varphi^{-1}$. We have \[\abs{\psi' - \id}_{s-2\sigma} \leq \abs{\varphi'^{-1} - \id}_{s-\sigma}\leq \frac{\abs{\varphi' - \id}_{s-\sigma}}{1 - \abs{\varphi' - \id}_{s-\sigma}}\leq \frac{2n}{2n -1}\frac{\abs{v}_s}{\sigma} \leq 2\frac{\abs{v}_s}{\sigma},\] by triangular and Cauchy inequalities.
	\end{itemize}
\end{proof}

\begin{cor}[Well definition of the operators $\phi$]
	\label{cor well def}For all $s,\sigma$
	\begin{itemize}[leftmargin=*]
		\item  if $g\in\G^{\sigma/n}_{s + \sigma}$, then $g^{-1}\in\A(\normal{T}^n_s,\normal{T}^n_{s+\sigma})$
		\item if $g\in\G^{\omega,\sigma^2/2n}_{s + \sigma}$, then $g^{-1}\in\A(\normal{T}^n_s,\normal{T}^n_{s+\sigma})$.
	\end{itemize}
	As a consequence, the operators $\phi$ in \eqref{normal operator}, \eqref{operator Herman} and \eqref{normal Russmann} are well defined.
\end{cor}
\begin{proof}
	We recall the form of $g\in\G^{\sigma/n}_{s+\sigma}$: \[g(\teta,r)= (\varphi(\teta), R_0(\teta) + R_1(\teta)\cdot r).\]
	$g^{-1}$ reads \[g^{-1}(\teta,r)=(\phi^{-1}(\teta), R_1^{-1}\circ\varphi^{-1}(\teta)\cdot(r - R_0\circ\varphi(\teta))).\]
	Up to rescaling norms by a factor $1/2$ like $\norm{x}_s := \frac{1}{2}\abs{x}$, the first statement is straightforward from theorem \ref{theorem well def}. By abuse of notations, we keep on indicating $\norm{x}_s$ with $\abs{x}_s$.\medskip\\
	Concerning those $g\in\G^{\omega,\sigma^2/2n}_{s+\sigma}$ we recall that $g^{-1}$ is given by
	\[g^{-1}(\teta,r)=(\varphi'^{-1}(\teta), \,^t\varphi'\circ\varphi^{-1}(\teta)\cdot r - \rho\circ\varphi^{-1}(\teta));\]
	if $\abs{\varphi^{-1} - \id}_{s}< \sigma$ and $\abs{\rho}_{s+\sigma}<\sigma/2$ with $\abs{r\cdot\varphi'\circ\varphi^{-1}(\teta)}_s < \sigma/2$ we get the wanted thesis. Just note that 
	\[\abs{\,^t(\varphi' - \id)\cdot r}_s \leq \frac{n\abs{r}}{\sigma}\abs{\varphi - \id}_{s+\sigma}\leq \sigma/2,\]
	the factor $n$ coming from the transposition. 
\end{proof}
\section{Estimates on the Lie brackets of vector fields}
This is just an adaptation to vector fields on $\normal{T}^n_{s+\sigma}$ of the analogous lemma for vector fields on the torus $\T^n_s$ in \cite{Poschel:2011}.
\begin{lemma}
	Let $f$ and $g$ be two real analytic vector fields on $\normal{T}^n_{s+\sigma}$. The following inequality holds
	\[\abs{\,\sq{f,g}\,}_s\leq \frac{2}{\sigma}\pa{1+\frac{1}{e}}\abs{f}_{s+\sigma}\abs{ g}_{s+\sigma}.\]
	\label{lemma lie brackets}
\end{lemma}
\begin{proof}
	Consider $f=(f^\teta,f^r) = \sum_{j=1}^n f^{\teta_j}\base{\teta_j} + f^{r_j}\base{r_j}$ and $g=(g^\teta,g^r) = \sum_{j=1}^n g^{\teta_j}\base{\teta_j} + g^{r_j}\base{r_j}$.
	From the definition of the Lie Brackets we have $[{f},{g}] = \sum_k f(g^k) - {g}(f^k),$ where every component $k$ reads
	\begin{align*}
		[f,g]^k &= \sum_{j=1}^n (f^{\teta_j}\frac{\partial g^k}{\partial\teta_j} + f^{r_j}\frac{\partial g^k}{\partial r_j}) - (g^{\teta_j}\frac{\partial f^k}{\partial \teta_j} + g^{r_j}\frac{\partial f^k}{\partial r_j})\\
		&= (Dg\cdot f - Df\cdot g)^k. 
	\end{align*}
	We observe that for an holomorphic function $h: \normal{T}^n_{s+\sigma}\to \C$, one has 
	\[\abs{\frac{\partial h}{\partial r_j}}_s = \sum_k \abs{\frac{\partial h _k(r)}{\partial r_j}}_s e^{\abs{k} s}\leq \sum_k \frac{1}{\sigma}{\abs{h_k(r)}_{s+\sigma}} e^{\abs{k}s}\leq  \frac{1}{\sigma}\abs{h}_{s+\sigma},\]
	and 
	\begin{align*}
		\abs{\frac{\partial h}{\partial \teta_j}}_s &= \sum_k \abs{k_j}\abs{h_k(r)}_s e^{\abs{k}s}\leq \sum_k \abs{k} \abs{h_k(r)}_s e^{\abs{k} (s+\sigma)}\, e^{-\abs{k}\sigma}\\
		& \leq \frac{1}{e\sigma}\sum_k \abs{h_k(r)}_{s+\sigma}e^{\abs{k}(s+\sigma)} = \frac{1}{e\sigma}\abs{h}_{s+\sigma},
	\end{align*}
	where we bound $\abs{k}e^{-\abs{k}\sigma}$ with the maximum attained by $x e^{-x\sigma},\, x > 0$, in $1/\sigma$, that is $1/e\sigma$. \\ Therefore, consider $f$ and $g$ in their Fourier's expansion, $Dg\cdot f$ read
	\[Dg \cdot f = \sum_{k,\ell}\,i {k\cdot f^\theta_\ell}g_k e^{i (k+\ell)\teta} + D_r g_k\cdot f^{r}_\ell\, e^{i (k + \ell)\cdot\teta} = \sum_{k,\ell} i\,k\cdot f^{\teta}_{\ell - k} g_k \,e^{i\ell\cdot\teta}+ D_r g_k\cdot f^{r}_{\ell-k}e^{i\ell\cdot\teta}. \]
	Passing to norms we have the following inequality
	\[\abs{Dg\cdot f}_s\leq \sum_{k,\ell} \abs{k}\abs{f^{\teta}_{\ell - k}}\abs{g_k} e^{\abs{k}s}e^{\abs{\ell-k}s} + \abs{D_r g_k}\abs{f^{r}_{\ell - k}} e^{\abs{k}s}e^{\abs{\ell -k }s}\leq\]
	\begin{align*}
		&\leq \sum_{k,\ell}\abs{k}\abs{g_k} e^{-\abs{k}\sigma} e^{\abs{k}(s+\sigma) }\abs{f^{\teta}_{\ell - k}}  e^{\abs{\ell-k}s} + \abs{D_r g_k}e^{\abs{k}s}\,\abs{f^r_{\ell - k}} e^{\abs{\ell - k}s}\\
		&\leq \frac{1}{e\sigma}\abs{g}_{s+\sigma}\abs{f}_{s+\sigma} + \frac{1}{\sigma}\abs{g}_{s+\sigma}\abs{f}_{s+\sigma},
	\end{align*}
	which follows from the previous remark. Hence the lemma.
\end{proof}

\bibliography{biblio_BJ}
\bibliographystyle{alpha}
%\bibliographystyle{plain}
\end{document}
